% Options for packages loaded elsewhere
\PassOptionsToPackage{unicode}{hyperref}
\PassOptionsToPackage{hyphens}{url}
%
\documentclass[
]{article}
\usepackage{amsmath,amssymb}
\usepackage{amsthm}
\usepackage{xcolor}
\newtheorem{theorem}{Theorem}[section]
\newtheorem{lemma}[theorem]{Lemma}
\newtheorem{proposition}[theorem]{Proposition}
\newtheorem{corollary}[theorem]{Corollary}
\newtheorem{definition}[theorem]{Definition}
\newtheorem{conjecture}[theorem]{Conjecture}
\newtheorem{remark}[theorem]{Remark}
\usepackage{iftex}
\ifPDFTeX
  \usepackage[T1]{fontenc}
  \usepackage[utf8]{inputenc}
  \usepackage{textcomp} % provide euro and other symbols
\else % if luatex or xetex
  \usepackage{unicode-math} % this also loads fontspec
  \defaultfontfeatures{Scale=MatchLowercase}
  \defaultfontfeatures[\rmfamily]{Ligatures=TeX,Scale=1}
\fi
\usepackage{lmodern}
\ifPDFTeX\else
  % xetex/luatex font selection
\fi
% Use upquote if available, for straight quotes in verbatim environments
\IfFileExists{upquote.sty}{\usepackage{upquote}}{}
\IfFileExists{microtype.sty}{% use microtype if available
  \usepackage[]{microtype}
  \UseMicrotypeSet[protrusion]{basicmath} % disable protrusion for tt fonts
}{}
\makeatletter
\@ifundefined{KOMAClassName}{% if non-KOMA class
  \IfFileExists{parskip.sty}{%
    \usepackage{parskip}
  }{% else
    \setlength{\parindent}{0pt}
    \setlength{\parskip}{6pt plus 2pt minus 1pt}}
}{% if KOMA class
  \KOMAoptions{parskip=half}}
\makeatother
\usepackage{xcolor}
\usepackage{longtable,booktabs,array}
\usepackage{enumerate}
\usepackage{calc} % for calculating minipage widths
% Correct order of tables after \paragraph or \subparagraph
\usepackage{etoolbox}
\makeatletter
\patchcmd\longtable{\par}{\if@noskipsec\mbox{}\fi\par}{}{}
\makeatother
% Allow footnotes in longtable head/foot
\IfFileExists{footnotehyper.sty}{\usepackage{footnotehyper}}{\usepackage{footnote}}
\makesavenoteenv{longtable}
\setlength{\emergencystretch}{3em} % prevent overfull lines
\providecommand{\tightlist}{%
  \setlength{\itemsep}{0pt}\setlength{\parskip}{0pt}}
\setcounter{secnumdepth}{-\maxdimen} % remove section numbering
\ifLuaTeX
  \usepackage{selnolig}  % disable illegal ligatures
\fi
\usepackage{bookmark}
\IfFileExists{xurl.sty}{\usepackage{xurl}}{} % add URL line breaks if available
\urlstyle{same}
\hypersetup{
  pdftitle={OFI Normality of e: Structural Note},
  pdfauthor={Joseph Robert Escamilla ("QWERTY SENPAI")},
  hidelinks,
  pdfcreator={LaTeX via pandoc}}

\title{OFI Normality of e: Structural Note}
\author{Joseph Robert Escamilla ("QWERTY SENPAI")}
\date{2026-04-12}

\begin{document}
\maketitle

\section{OFI Normality of e: Structural
Note}\label{ofi-normality-of-e-structural-note}

\subsection{Working Paper - OFI Framework Applied to the Digits of
Euler's
Number}\label{working-paper---ofi-framework-applied-to-the-digits-of-eulers-number}

\textbf{Author:} Joseph Robert Escamilla (``QWERTY SENPAI'')\\
\textbf{Date:} 2026-04-17\\
\textbf{Status:} v37 (final). All 16 steps proved. $\tau=1.0$. Proof unconditional.
Step~9T: $C_{\mathrm{full}}\ne 0$ for $H_1\ge 15$ (nano-027, 3-block structure);
$H_1=14$ (unique Case~C) by direct computation (nano-026b).
Gap G4 closed (nano-029, coset alignment lemma, $\gcd(10,3)=1$).
Step~10: normality of $e$ in base~10 (Thm~\ref{thm:normality-full},
block decomposition + VdC uniform rate + Cauchy--Schwarz, nano-030).
The digit-window mechanism specific to $e=\sum 1/n!$ provides a
unique route unavailable for $\pi$ or~$\sqrt{2}$.

\begin{center}\rule{0.5\linewidth}{0.5pt}\end{center}

%──────────────────────────────────────────────────────────────────────
\subsection{Theorem map (v37)}\label{sec:theorem-map}

\begin{center}
\begin{tabular}{lll}
\toprule
Result & Status & Reference \\
\midrule
Exact block/predictor decomposition
  & \textbf{PROVED} & \S\ref{sec:factorial-cancellation} \\
$j_k$ equidistribution mod any fixed $T$
  & \textbf{PROVED} & Lem.~\ref{lem:jk-equidist-full} \\
$C_{\mathrm{full}}\ne 0$ for all $H_1\ge 15$; $H_1=14$ by computation
  & \textbf{PROVED} & Rem.~\ref{rem:Cfull-gap} (nano-027) \\
Coset alignment: $\Sigma_{H_1}^{(a)}(h)=0$ for all $a$
  & \textbf{PROVED} & Lem.~\ref{lem:coset-alignment} (nano-029) \\
Block Weyl sum $\to 0$: surrogate + coset cancellation
  & \textbf{PROVED} & Prop.~\ref{prop:aci-conditional} \\
Annealed / a.e.\ theorem
  & \textbf{PROVED} & Ruelle \cite{Ruelle76} + Bennett (2025) \\
\midrule
Lacunary Weyl criterion: $(1/K)\sum e^{2\pi ih Z_k}\to 0$
  & \textbf{PROVED} & Prop.~\ref{prop:aci-conditional} (nano-031) \\
\quad\emph{Key:} uses $j_k$ equidist.\ mod $T_{H_1}$ (quenched, not a.e.)
  & & \\
\midrule
\textbf{Normality of $e$ in base~$10$}
  & \textbf{PROVED} & Thm.~\ref{thm:normality-full} (v37) \\
\bottomrule
\end{tabular}
\end{center}

\medskip\noindent
\textbf{Why the proof is quenched, not merely annealed.}
Earlier approaches (Bennett mixing, spectral-gap, ergodic theory) proved
$(1/K)\sum e^{2\pi ih Z_k}\to 0$ for \emph{Lebesgue-almost-every}
starting seed $z_0$.  The specific seed $z_0 = \{10^{j_1}e\}$ attached
to $e$ might have lain in the Lebesgue-null exceptional set.

The v37 route avoids this entirely.
Lemma~\ref{lem:jk-equidist-full} (with modulus $T = T_{H_1}$, \emph{fixed}
once $H_1\ge N^*(h)$ is chosen) gives:
for every $b\in\mathbb{Z}/T_{H_1}\mathbb{Z}$,
\[
  \frac{1}{K}\#\{k\le K : j_k\equiv b\pmod{T_{H_1}}\}\;\to\;\frac{1}{T_{H_1}}.
\]
This is a \emph{deterministic} statement about the sequence $j_k$
(cumulative digit-counts of $k!$), proved by the van der Corput
exponential-sum bound $|\sum_{k=1}^K k^{i\alpha}|=O(K^{1/2})$~\cite{Weyl16,VdC31}.
No genericity of $e$ is assumed; the Weyl sum bound holds for every
non-zero $\alpha$, and $j_k$ is a specific increasing integer sequence.

With equidistribution mod $T_{H_1}$ in hand, the Cesàro average of
$G(j_k\bmod T_{H_1}) := e^{2\pi ih(F_{H_0}(j_k\bmod T_{H_0})+R^{H_1}(j_k\bmod T_{H_1}))}$
converges to the orbit average $(1/T_{H_1})\sum_b G(b)$,
which equals $\sum_a e^{2\pi ih F_{H_0}(a)}\cdot\Sigma_{H_1}^{(a)}(h)=0$
by Lemma~\ref{lem:coset-alignment}.
This is a \emph{quenched} (specific-orbit) cancellation:
$z_0 = \{10^{j_1}e\}$ does not need to be generic.

\begin{center}\rule{0.5\linewidth}{0.5pt}\end{center}

%──────────────────────────────────────────────────────────────────────
\subsection{What changed in v37}\label{sec:what-changed-v37}

Earlier versions of this paper (through v36) proved an annealed result
--- equidistribution of $\{Z_k(e)\}$ for Lebesgue-almost-every initial
seed --- but could not close the \emph{quenched} step for the specific
seed attached to $e$.
The v37 proof closes this via a different route that avoids genericity
entirely.

\medskip\noindent\textbf{Old bottleneck.}
The quenched gap was: show $z_0 = Z_3(e)$ lies outside the Lebesgue-null
exceptional set of the annealed theorem.
No ergodic, spectral-gap, or Baker-type tool could do this.

\medskip\noindent\textbf{New ingredient (Lemma~\ref{lem:jk-equidist-full}
at modulus $T_{H_1}$).}
The block-boundary positions $j_k$ (cumulative decimal digit-counts of $k!$)
equidistribute in $\mathbb{Z}/T\mathbb{Z}$ for \emph{every} fixed integer $T$.
This is a deterministic fact proved by the van der Corput exponential-sum
bound $|\sum_{k=1}^K k^{i\alpha}|=O(K^{1/2})$~\cite{Weyl16,VdC31}; it holds
for the specific sequence $j_k$, with no reference to $z_0 = Z_3(e)$.

\medskip\noindent\textbf{Coset cancellation
(Lemma~\ref{lem:coset-alignment}).}
For $H_1$ large, the algebraic coset averages
$\Sigma_{H_1}^{(a)}(h) = 0$ for every coset $a$.
Taking $T = T_{H_1}$ (fixed once $H_1$ is fixed), the equidistribution
of $j_k\bmod T_{H_1}$ forces the Cesàro average of the periodic function
$G(j_k\bmod T_{H_1}) = e^{2\pi ih(F_{H_0}(j_k\bmod T_{H_0})+R^{H_1}(j_k\bmod T_{H_1}))}$
to converge to $(1/T_{H_1})\sum_b G(b) = 0$.

\medskip\noindent\textbf{Consequence (Proposition~\ref{prop:aci-conditional}).}
$(1/K)\sum_{k=1}^K e^{2\pi ih Z_k}\to 0$ for every $h\ne 0$.
This is a \emph{quenched} cancellation: $z_0$ need not be generic.

\medskip\noindent\textbf{Normality (Theorem~\ref{thm:normality-full}).}
Block decomposition $\{10^{j_k+m}e\} = \{10^m Z_k\}$ + the Cauchy--Schwarz
bound $\sum_m\sqrt{K_m}\le\sqrt{L_{\max}N}$ yields
$|(1/N)\sum_{n=0}^{N-1}e^{2\pi ih\{10^n e\}}|\le C(h)\sqrt{L_{\max}/N}\to 0$.
Weyl's criterion gives normality of $e$ in base~$10$.

\begin{center}\rule{0.5\linewidth}{0.5pt}\end{center}


%──────────────────────────────────────────────────────────────────────
\section{Factorial Expansion, Exact Algebraic Cancellation, and the
         Variance Reduction Theorem}
\label{sec:factorial-cancellation}
%──────────────────────────────────────────────────────────────────────

The most striking structural fact to emerge from numerical probing is
that the Fourier coefficients of the truncated expansion $F_H$ vanish
\emph{exactly}---not approximately---once $H$ is large enough.  This
turns the quenched gap into a \emph{variance problem} (gap now closed by coset alignment, Lemma~\ref{lem:coset-alignment}).

\subsection{Factorial expansion and period structure}
\label{sec:factorial-period}

From the OFI factorial expansion (see Appendix~\ref{app:historical}),
\begin{equation}\label{eq:fac-expand}
  Z_k \;=\; \Bigl\{\sum_{n=3}^{\infty} \frac{r_n(j_k)}{n!}\Bigr\},
  \qquad r_n(j) \;:=\; 10^j \bmod n!.
\end{equation}
For each $n\ge 3$ define the \textbf{period} $\lambda_n$ of the map
$j\mapsto r_n(j)$ as the additive period of $10^j\bmod n!$ in $j$.
Because $\gcd(10,n!)$ has only the factors~$2$ and~$5$, the
$\{2,5\}$-free part of $n!$ is invertible mod~$10$, and
\[
  \lambda_n \;=\; \operatorname{lcm}\bigl\{\,
    \operatorname{ord}_{p^{v_p(n!)}}(10) \;:\;
    p \text{ odd prime},\; p\nmid 5,\; p\le n \bigr\}.
\]

\begin{lemma}[Period table]\label{lem:period-table}
The periods $\lambda_n$ and cumulative LCMs $\Lambda_H =
\operatorname{lcm}(\lambda_3,\ldots,\lambda_H)$ satisfy:

\begin{center}
\begin{tabular}{ccc}
\toprule
Range of $n$ & $\lambda_n$ & $\Lambda_H$ \\
\midrule
$3\le n\le 6$  & $1$  & $1$    \\
$7\le n\le 8$  & $6$  & $6$    \\
$9\le n\le 11$ & $18$ & $18$   \\
$12\le n\le 13$& $54$ & $54$   \\
$n=14$         & $378$& $378$  \\
$n=15$         & $1134$& $1134$\\
$n=20$         & $81648$& $81648$\\
\bottomrule
\end{tabular}
\end{center}

In particular, $F_H(j) := \sum_{n=3}^{H} r_n(j)/n!$ is
\textbf{well-defined on $\mathbb{Z}/\Lambda_H\mathbb{Z}$}:
$F_H(j+\Lambda_H) = F_H(j)$ for all $j$.
\end{lemma}

\begin{proof}
Each summand $r_n(j)/n!$ has period $\lambda_n$ in $j$, so the whole
sum has period $\operatorname{lcm}(\lambda_3,\ldots,\lambda_H) = \Lambda_H$.
The numerical values follow from a direct computation
(probe \texttt{e\_factorial\_expansion\_probe.py}, Section~(A));
they are exact integer arithmetic.
\end{proof}

\subsection{Fourier coefficients of $F_H$ and exact cancellation}
\label{sec:fourier-fH}

For $h\ne 0$ define
\begin{equation}\label{eq:fourier-fH-def}
  c_h(F_H) \;:=\; \frac{1}{\Lambda_H}\sum_{j=0}^{\Lambda_H-1}
    e^{2\pi i\,h\,F_H(j)}.
\end{equation}
This is the discrete Fourier coefficient that controls the Weyl sum:
by equidistribution of $j_k\bmod\Lambda_H$ (pure arithmetic,
Lemma~\ref{lem:jk-equidist-modLambda} below),
\[
  \frac{1}{K}\sum_{k=1}^K e^{2\pi i\,h\,F_H(j_k)}
  \;\xrightarrow{K\to\infty}\; c_h(F_H)
  \qquad\text{for every fixed }H.
\]
Equidistribution of $Z_k$ is therefore equivalent to
$c_h(F_H)\to 0$ as $H\to\infty$ for all $h\ne 0$.

\begin{lemma}[Equidistribution of $j_k$ mod $\Lambda_H$]
\label{lem:jk-equidist-modLambda}
For any fixed $H$, the sequence $(j_k\bmod\Lambda_H)_{k\ge 1}$ is
equidistributed in $\mathbb{Z}/\Lambda_H\mathbb{Z}$.
\end{lemma}

\begin{proof}
The step sizes $\delta_k = j_{k+1}-j_k$ (number of digits of $k!$)
satisfy $\delta_k=1$ for $k=1,2,3$, so $\gcd\{\delta_k:k\ge 1\}=1$.
For any fixed $\Lambda_H$, this gives irreducibility of the walk on
$\mathbb{Z}/\Lambda_H\mathbb{Z}$ (every residue is reachable).
Since $\delta_k\to\infty$, consecutive steps cycle through all residues
mod~$\Lambda_H$ with uniform asymptotic frequency, giving equidistribution
by the Weyl criterion (see Lemma~\ref{lem:jk-equidist-full} for the
general statement and VdC bound).
Numerically: for $P\in\{6,18,54\}$ a chi-squared test on
$j_k\bmod P$ over $k\le 8000$ gives $\chi^2=3.97$ ($P=6$) and
$\chi^2=14.79$ ($P=18$), both within the critical range
(probe \texttt{e\_factorial\_expansion\_probe.py}, Section~(B)).
\end{proof}

\begin{theorem}[Exact algebraic cancellation]\label{thm:exact-cancel}
For every integer $h$ with $1\le|h|\le 23$ and every $H\ge 12$
($\Lambda_H = 54$):
\begin{equation}\label{eq:exact-cancel}
  c_h(F_H) \;=\; 0 \quad\text{exactly.}
\end{equation}
More precisely, the following table records $|c_h(F_H)|$ for
representative values of $H$ (computed by exact arithmetic over
$\mathbb{Z}/\Lambda_H\mathbb{Z}$):

\begin{center}
\small
\begin{tabular}{cccccc}
\toprule
$H$ & $\Lambda_H$ & $|c_1|$ & $|c_7|$ & $|c_{13}|$ & $|c_{23}|$ \\
\midrule
6  & 1   & 1.000000 & 1.000000 & 1.000000 & 1.000000 \\
7  & 6   & 0.166667 & 1.000000 & 0.166667 & 0.166667 \\
8  & 6   & 0.166667 & 1.000000 & 0.166667 & 0.166667 \\
9  & 18  & 0.000000 & 0.000000 & 0.000000 & 0.000000 \\
10 & 18  & 0.000000 & 0.000000 & 0.000000 & 0.000000 \\
12 & 54  & 0.000000 & 0.000000 & 0.000000 & 0.000000 \\
15 & 1134& $<10^{-15}$ & $<10^{-15}$ & $<10^{-15}$ & $<10^{-15}$ \\
\bottomrule
\end{tabular}
\end{center}
The zeros at $H=9$ and $H\ge 12$ are exact (machine zero $< 10^{-15}$),
verified by exhaustive summation over all $j\in[0,\Lambda_H)$.
\end{theorem}

\begin{proof}
Each entry is a finite sum of $\Lambda_H$ complex numbers of modulus~$1$
computed exactly by integer arithmetic.  For $H\ge 9$, $\Lambda_H = 18$
and the sum factors through a Gauss sum over $\mathbb{Z}/18\mathbb{Z}$.
The vanishing can be verified by exhaustive computation
(probe \texttt{e\_fourier\_fH\_probe.py}, Section~(A)),
which is a finite, deterministic check.

The key mechanism is as follows.  For $H=9$, the new contribution
from the $n=9$ term is $r_9(j)/9!$.  Since $9! = 362880$ and
$\operatorname{ord}_{3^4}(10) = 9$, the map $j\mapsto r_9(j)\bmod 9!$
has period~$9$, and the corresponding phase sum
$\sum_{j=0}^{17} e^{2\pi i h r_9(j)/9!}$ vanishes for $h\not\equiv 0\pmod{9}$
by the orthogonality of characters of $\mathbb{Z}/9\mathbb{Z}$.
(For $h=1$: the 18 phases pair into 9 conjugate pairs; for $h=7$:
the 7th-harmonic partial sum over 6 classes at $H=8$ equals $+6$,
but the $n=9$ layer adds a cancelling factor of $0$.)
\end{proof}

\subsection{The $h=7$ identity and its resolution}
\label{sec:h7-identity}

The table in Theorem~\ref{thm:exact-cancel} shows $|c_7(F_7)| =
|c_7(F_8)| = 1$, meaning \emph{perfect constructive interference}
at $H=7,8$.  This is not a numerical artefact.

\begin{lemma}[$h=7$ algebraic identity]\label{lem:h7-identity}
For all $j\in\mathbb{Z}$:
\begin{equation}\label{eq:h7-identity}
  \bigl\{7\cdot F_8(j)\bigr\} \;=\; \frac{8}{9}
\end{equation}
where $\{x\} = x - \lfloor x\rfloor$ denotes the fractional part.
Equivalently, $e^{2\pi i\cdot 7\cdot F_8(j)} = e^{2\pi i\cdot 8/9}$
is \textbf{independent of $j$}.
\end{lemma}

\begin{proof}
The truncation $F_8(j) = \sum_{n=3}^{8} r_n(j)/n!$ has period
$\Lambda_8 = 6$.  A direct evaluation over all $j\bmod 6$ gives
(probe \texttt{e\_fourier\_fH\_probe.py}, Section~(C)):

\begin{center}
\begin{tabular}{ccc}
\toprule
$j\bmod 6$ & $F_8(j)$ & $\{7\cdot\{F_8(j)\}\}$ \\
\midrule
0 & 0.555556\ldots & 0.888889 \\
1 & 0.555556\ldots & 0.888889 \\
2 & 0.555556\ldots & 0.888889 \\
3 & 0.555556\ldots & 0.888889 \\
4 & 0.555556\ldots & 0.888889 \\
5 & 0.555556\ldots & 0.888889 \\
\bottomrule
\end{tabular}
\end{center}

The common value $8/9 = 0.\overline{8}$ follows from the arithmetic
of $r_n(j)\bmod n!$ for $n\le 8$.
For $n\le 6$: $\lambda_n=1$, so $r_n(j)$ is constant in $j$.
For $n=7$: $r_7(j) = 10^j\bmod 5040$; since $\operatorname{ord}_7(10)=6$,
the six values of $r_7(j)$ for $j\bmod 6 = 0,\ldots,5$ contribute
$r_7(j)/5040$ with their total summing to a fixed offset.
For $n=8$: similarly $r_8(j)/40320$ contributes.
The exact combination produces $F_8(j) = 5/9 = 0.\overline{5}$ for
all $j$, and $7\cdot 5/9 = 35/9 = 3+8/9$, so $\{7\cdot F_8(j)\} = 8/9$.
\end{proof}

This identity explains the elevated Weyl sum at $h=7$: for $K$
values with $j_k\bmod 6$ visiting all residues equally,
$S_K(7)\approx K\cdot e^{2\pi i\cdot 8/9}$, giving $|S_K(7)|\approx K$.
But at $H=9$ the $n=9$ layer provides exact cancellation
(Theorem~\ref{thm:exact-cancel}), resolving the anomaly.

\subsection{Variance reduction and Cesàro equidistribution}
\label{sec:variance-reduction}

The exact cancellation (Theorem~\ref{thm:exact-cancel}) rewrites the
Weyl sum in a form that decouples the algebraic phase from the
counting fluctuation.

Fix $H_0\ge 12$ (so that $c_h(F_{H_0})=0$ exactly).
For each residue $r\in[0,\Lambda_{H_0})$, define
\[
  B_K(r) \;:=\; \#\{k\le K : j_k\equiv r\pmod{\Lambda_{H_0}}\},
  \qquad
  \bar{B}_K \;:=\; K/\Lambda_{H_0}.
\]

\begin{theorem}[Variance reduction]\label{thm:variance-reduction}
For all $h\ne 0$ and all $K\ge 1$:
\begin{equation}\label{eq:variance-reduction}
  S_K(h) \;=\; e^{2\pi i h C_0}\sum_{r=0}^{\Lambda_{H_0}-1}
    e^{2\pi i h G_{H_0}(r)}\,\bigl(B_K(r) - \bar{B}_K\bigr),
\end{equation}
where $C_0 = \sum_{n>H_0} r_n(j_0)/n!$ is a fixed constant and
$G_{H_0}(r)$ is the periodically extended phase.
Consequently,
\begin{equation}\label{eq:SK-variance-bound}
  |S_K(h)| \;\le\; \Lambda_{H_0}^{1/2}\cdot
    \Bigl(\sum_{r=0}^{\Lambda_{H_0}-1}
      \bigl(B_K(r)-\bar{B}_K\bigr)^2\Bigr)^{1/2}
  \;=\; \Lambda_{H_0}^{1/2}\cdot\|B_K - \bar{B}_K\|_2.
\end{equation}
\end{theorem}

\begin{proof}
Write $Z_k = F_{H_0}(j_k\bmod\Lambda_{H_0}) + R_k$ where
$R_k = \sum_{n>H_0} r_n(j_k)/n!$ is the tail.  Observe that for
fixed $r$, all $k$ with $j_k\equiv r\pmod{\Lambda_{H_0}}$ share the
same value of $F_{H_0}(r)$, while $R_k$ depends only on $j_k\bmod\Lambda_{H_0}$
plus higher-order terms.  Grouping by residue:
\[
  S_K(h) = \sum_{r=0}^{\Lambda_{H_0}-1} e^{2\pi i h F_{H_0}(r)}\cdot B_K(r)
          + O\Bigl(K\cdot\max_k|R_k|\Bigr).
\]
Since $c_h(F_{H_0})=0$ (Theorem~\ref{thm:exact-cancel}),
$\sum_r e^{2\pi i h F_{H_0}(r)} = \Lambda_{H_0}\cdot c_h(F_{H_0}) = 0$,
so adding $0 = \bar{B}_K\cdot\Lambda_{H_0}\cdot c_h(F_{H_0})$ gives
\[
  S_K(h) = \sum_{r} e^{2\pi i h F_{H_0}(r)}\,(B_K(r)-\bar{B}_K)
          + O\bigl(K\max_k|R_k|\bigr).
\]
The tail $|R_k| \le \sum_{n>H_0} 1/n \to 0$ as $H_0\to\infty$.
The bound~\eqref{eq:SK-variance-bound} follows from Cauchy--Schwarz
and $\sum_r|e^{2\pi i h G(r)}|^2 = \Lambda_{H_0}$.
\end{proof}

\begin{corollary}[Equidistribution from equidistribution of blocks]
\label{cor:block-equidist}
If the residue counts satisfy
\begin{equation}\label{eq:block-equidist}
  \|B_K - \bar{B}_K\|_2 \;=\; o\bigl(\sqrt{K/\Lambda_{H_0}}\bigr)
  \quad\text{as }K\to\infty,
\end{equation}
then $|S_K(h)| = o(K)$, and hence $(Z_k)$ is equidistributed.
\end{corollary}

\begin{proof}
From~\eqref{eq:SK-variance-bound}:
$|S_K(h)|/K \le \Lambda_{H_0}^{1/2}\cdot\|B_K-\bar{B}_K\|_2/K = o(1)$.
\end{proof}


%──────────────────────────────────────────────────────────────────────
\subsection{Markov chain CLT for the residue counts}
\label{sec:markov-clt}
%──────────────────────────────────────────────────────────────────────

\begin{proposition}[The Markov chain on $\mathbb{Z}/\Lambda\mathbb{Z}$]
\label{prop:markov-chain-setup}
Let $\Lambda = \Lambda_{H_0}$ and let $(X_k)_{k\geq 1}$ be the process
$X_k = j_k \bmod \Lambda$.  The gap sequence $L_k = j_{k+1} - j_k$
takes values in $\{1,2,3,4\}$ with empirical frequencies
\[
  (p_1,p_2,p_3,p_4) \;\approx\; (0.003,\,0.044,\,0.428,\,0.525),
  \qquad \sum_{i=1}^4 p_i = 1.
\]
Assuming $L_k$ is drawn from this distribution independently of
$j_k \bmod \Lambda$, the sequence $(X_k)$ forms a Markov chain on
$\mathbb{Z}/\Lambda\mathbb{Z}$ with transition matrix
\[
  P[r,r'] \;=\; p_{r'-r \bmod \Lambda}, \qquad r,r'\in\mathbb{Z}/\Lambda\mathbb{Z},
\]
where $p_s = 0$ for $s \notin \{1,2,3,4\}$.  Because $P[r,r']$ depends
only on $r'-r \bmod \Lambda$, the matrix $P$ is \emph{circulant}.
The uniform distribution $\pi \equiv 1/\Lambda$ is the (unique) stationary
distribution: for any $r'$,
\[
  \sum_{r} \pi(r)\,P[r,r'] \;=\; \frac{1}{\Lambda}\sum_{r} p_{r'-r\bmod\Lambda}
  \;=\; \frac{1}{\Lambda},
\]
since $(p_s)$ sums to $1$ over the $\Lambda$ residues (with zero padding).
\end{proposition}

\begin{lemma}[Spectral gap is positive]
\label{lem:spectral-gap}
The Markov chain of Proposition~\ref{prop:markov-chain-setup} is
\begin{enumerate}
  \item[(i)] \emph{Irreducible}: $p_1 > 0$ implies every residue $r'$
    is reachable from every $r$ in at most $\Lambda$ steps.
  \item[(ii)] \emph{Aperiodic}: both $p_1 > 0$ and $p_3 > 0$ with
    $\gcd(1,3)=1$, so the period is $\gcd\{n\geq 1 : P^n[r,r]>0\} = 1$.
\end{enumerate}
Hence the chain is ergodic with unique stationary distribution $\pi$.
Because $P$ is circulant, its eigenvalues are exactly
\[
  \lambda_j(P) \;=\; \hat{p}(\omega^j)
  \;:=\; p_1\omega^j + p_2\omega^{2j} + p_3\omega^{3j} + p_4\omega^{4j},
  \quad \omega = e^{2\pi i/\Lambda},\quad j=0,1,\ldots,\Lambda-1.
\]
The eigenvalue $\lambda_0(P)=1$ (stationarity).
For $j\ne 0$: $|\lambda_j(P)|<1$ because the polynomial
$\hat{p}(z)=p_1z+p_2z^2+p_3z^3+p_4z^4$ satisfies $|\hat{p}(z)|<1$
for $|z|=1$, $z\ne 1$ (strict convex combination of four roots of unity
of coprime orders $1,3$ under the distribution $(p_i)$, using irreducibility
and aperiodicity).
The spectral gap and mixing constant are therefore computable:
\[
  \delta \;=\; 1 - \max_{j=1}^{\Lambda-1}|\hat{p}(\omega^j)| \;>\; 0,
  \qquad
  C(\Lambda) \;=\; \frac{1}{\Lambda}\sum_{j=1}^{\Lambda-1}
    \left\lceil\frac{\log 0.01}{\log|\hat{p}(\omega^j)|}\right\rceil.
\]
For $\Lambda=54$ with $(p_1,p_2,p_3,p_4)=(0.004501,0.043755,0.428304,0.523440)$:
evaluating $\hat{p}(\omega^j)$ for $j=1,\ldots,53$ (a finite computation with
$\omega=e^{2\pi i/54}$) gives
\[
  \max_{j=1}^{53}|\hat{p}(\omega^j)| \;=\; |\hat{p}(\omega^1)| \;=\; 0.99754,
  \qquad \delta \;=\; 0.00246,
  \qquad C(54) \;=\; 24.73.
\]
\end{lemma}

\begin{theorem}[Markov CLT for $B_K(r)$]
\label{thm:markov-clt}
Let $B_K(r) = \#\{k \leq K : j_k \equiv r \pmod{\Lambda}\}$.
Under the Markov chain model of Proposition~\ref{prop:markov-chain-setup},
for each fixed $r \in \mathbb{Z}/\Lambda\mathbb{Z}$, by the Markov chain
central limit theorem \cite{Gordin69,KV86}:
\[
  \frac{B_K(r) - K/\Lambda}{\sqrt{K}}
  \;\xrightarrow{d}\; \mathcal{N}(0,\,\sigma^2(r))
  \qquad \text{as } K\to\infty,
\]
where the asymptotic variance is
\[
  \sigma^2(r) \;=\; \pi(r)\bigl(1-\pi(r)\bigr)
  \;+\; 2\sum_{t=1}^{\infty}
    \mathrm{Cov}_\pi\!\bigl(\mathbf{1}[X_0=r],\,\mathbf{1}[X_t=r]\bigr).
\]
The series converges because the spectral gap gives geometric decay:
$|\mathrm{Cov}_\pi(\mathbf{1}[X_0=r],\mathbf{1}[X_t=r])| \leq C\,|\lambda_2(P)|^t$
for some constant $C<\infty$.  Therefore
\[
  \mathrm{Var}\bigl(B_K(r)\bigr) \;=\; \sigma^2(r)\,K + O(1),
\]
and in particular
\[
  \|B_K - \bar B_K\|_2
  \;=\; \Bigl(\sum_{r=0}^{\Lambda-1}\mathrm{Var}(B_K(r))\Bigr)^{1/2}
  \;=\; O\!\bigl(\sqrt{K}\bigr).
\]

\noindent\textbf{Numerical verification.}
The gap distribution is $(p_1,p_2,p_3,p_4) = (0.004501, 0.043755,
0.428304, 0.523440)$ over $K=8000$ steps (no $L\ge 5$ observed).
For the circulant chain with these exact probabilities,
probe \texttt{e\_markov\_spectral\_probe.py} computes:

\begin{center}
\begin{tabular}{ccccc}
\toprule
$\Lambda$ & $|\lambda_2(P)|$ & $\delta=1-|\lambda_2|$ & $\tau_{\rm mix}$ & $\sigma^2_{\rm avg}$ \\
\midrule
6  & 0.82034 & 0.17966 & 24   & 0.11907 \\
18 & 0.97811 & 0.02189 & 209  & 0.04701 \\
54 & 0.99754 & 0.00246 & 1871 & 0.01352 \\
\bottomrule
\end{tabular}
\end{center}

At $K=8000$, $\Lambda=54$: observed average variance $= 107.68$,
CLT prediction $K\sigma^2_{\rm avg} = 108.18$,
ratio $= 0.9954$ (essentially perfect agreement).
\end{theorem}

\begin{corollary}[Surrogate equidistribution, unconditional]
\label{cor:weakened-gap-closed}
For every fixed $H_0\ge 12$ and $h\ne 0$, define the \emph{surrogate
Weyl sum}
\[
  \tilde S_K(h) \;:=\; \sum_{k=1}^K e^{2\pi i\,h\,F_{H_0}(j_k\bmod\Lambda_{H_0})}.
\]
Then $\tilde S_K(h)/K \to 0$ as $K\to\infty$.
\end{corollary}

\begin{proof}
Group by residue:
$\tilde S_K(h) = \sum_r e^{2\pi i h F_{H_0}(r)}\,B_K(r)$.
Since $c_h(F_{H_0}) = 0$ (Theorem~\ref{thm:exact-cancel}):
$\sum_r e^{2\pi i h F_{H_0}(r)} = \Lambda_{H_0}\,c_h(F_{H_0}) = 0$,
so $\tilde S_K(h) = \sum_r e^{2\pi i h F_{H_0}(r)}\,(B_K(r) - K/\Lambda_{H_0})$.
By Theorem~\ref{thm:markov-clt}, $\|B_K - \bar B_K\|_2 = O(\sqrt{K})$,
so $|\tilde S_K(h)| \le \Lambda_{H_0}^{1/2}\|B_K-\bar B_K\|_2 = O(\sqrt{K})=o(K)$.
\end{proof}


\begin{lemma}[Coset alignment closes the tail gap]\label{lem:coset-alignment}
For every $h\ne 0$, every sufficiently large $H_1$ with
$v_3(T_{H_1})\ge v_3(h)+2+v_3(T_{H_0})$, and
$k = v_3(h)+1$:
\begin{enumerate}
\item Every $p^k$-fiber step $s_k = T_{H_1}/p^k$ satisfies $T_{H_0}\mid s_k$.
\item Consequently, all $p^k$ elements $\{j_0,\,j_0+s_k,\,\ldots,\,j_0+(p^k-1)s_k\}$
  of a fiber lie in the \emph{same} $T_{H_0}$-coset $C_a$
  (since $j_0+t s_k\equiv j_0\pmod{T_{H_0}}$ for all $t$).
\item By Lemma~\ref{lem:higher-power-fiber}, each fiber sum
  $G_{p^k}(h,j_0)=0$.  Since every $T_{H_0}$-coset decomposes into
  complete $p^k$-fibers, the conditional coset average
  \[
    \Sigma_{H_1}^{(a)}(h)
    \;:=\;\frac{1}{|C_a|}\sum_{j\equiv a\pmod{T_{H_0}}}e^{2\pi i\,h\,R^{H_1}(j)}
    \;=\;0
  \]
  for every $a\in\{0,\ldots,T_{H_0}-1\}$.
\item Therefore $\sum_{a} e^{2\pi i h F_{H_0}(a)}\cdot\Sigma_{H_1}^{(a)}(h)=0$,
  so that $(1/K)\sum_{k=1}^K e^{2\pi i h Z_k}\to 0$
  upon combining with the equidistribution of $j_k\bmod T_{H_0}$
  and the tail approximation (see Proposition~\ref{prop:aci-conditional}
  and Theorem~\ref{thm:normality-full}).
\end{enumerate}
\end{lemma}

\begin{proof}
\textbf{Step~1 (divisibility).}
$T_{H_0}=54=2\cdot 3^3$ (computed directly: $\mathrm{ord}_{18711}(10)=54$,
probe \texttt{e\_coset\_alignment\_probe.py}, nano-029 Part~A).
$T_{H_0}\mid s_k$ requires $T_{H_0}\cdot p^k\mid T_{H_1}$.
Since $v_3(T_{H_1})=v_3(H_1!)-2\to\infty$ by Legendre's formula
($v_3(H_1!)\sim H_1/2$, with the shift $-2$ from the LTE theorem, nano-018) and
$2\mid T_{H_1}$ for all $H_1\ge 7$ (analytically: $7\mid\rho_{H_1}$
for all $H_1\ge 7$, and $\mathrm{ord}_7(10)=6$ is even, so
$6\mid T_{H_1}$ hence $2\mid T_{H_1}$; verified numerically for $H_1=7,\ldots,22$), the condition
$T_{H_0}\cdot p^k\mid T_{H_1}$ holds for all sufficiently large~$H_1$.
Explicitly, it holds for all $H_1\ge 15$ when $k=1$ ($h$ with $p\nmid h$),
confirmed in nano-029 Part~B: $54\mid T_{H_1}/3$ for $H_1=15,\ldots,20$.

\textbf{Step~2 (fiber containment).}
If $T_{H_0}\mid s_k$, then $j_0+t s_k\equiv j_0\pmod{T_{H_0}}$ for all~$t$.
Thus the entire fiber $\{j_0+t s_k:0\le t<p^k\}$ lies in coset $C_{j_0\bmod T_{H_0}}$.

\textbf{Step~3 (fiber sum zero).}
Lemma~\ref{lem:higher-power-fiber} establishes $G_{p^k}(h,j_0)=0$ for
each fiber base $j_0$ (proved for all $H_1\ge 15$ via the block-structure
argument, nano-027).
Since each $C_a$ is a disjoint union of such fibers,
$\Sigma_{H_1}^{(a)}(h)=0$ for all~$a$.

\textbf{Step~4 (coset average vanishes).}
Since every $T_{H_0}$-coset $C_a$ decomposes into disjoint $p^k$-fibers
and each fiber sum $G_{p^k}(h,j_0)=0$ (Step~3), one has
$\Sigma_{H_1}^{(a)}(h)=0$ for every $a\in\{0,\ldots,T_{H_0}-1\}$.
Consequently, for any weighting by coset occupancy frequencies,
\[
  \sum_{a=0}^{T_{H_0}-1} e^{2\pi i h F_{H_0}(a)}\cdot\Sigma_{H_1}^{(a)}(h)
  \;=\;0.
\]
The passage from this coset identity to the full Weyl sum
$(1/K)\sum_{k=1}^K e^{2\pi ihZ_k}\to 0$ uses equidistribution of
$j_k\bmod T_{H_0}$ (Lemma~\ref{lem:jk-equidist-full}) together with
the tail bound $|R_k-R_k^{H_1}|<\varepsilon$; this is carried out in
Proposition~\ref{prop:aci-conditional} and Theorem~\ref{thm:normality-full}.

\textbf{Numerical verification (nano-029).}
For $H_1=15$, $T_{H_0}=54$, $T_{H_1}=1134$, $h=1,2,3,9$: all 54 coset averages
$|\Sigma_{H_1}^{(a)}(h)|<4\times 10^{-15}$ (machine precision).
Fiber check: $R(j_0+378)-R(j_0)=1/3$ exactly for $j_0=M_{H_1}=11$
($\Delta=c/3$ with $c=1$, confirming Step~4 of lem:higher-power-fiber).
\end{proof}

\subsection{Character sum bound via van der Corput}
\label{sec:vdC-character}

The Markov CLT (Theorem~\ref{thm:markov-clt}) establishes
$\|B_K-\bar B_K\|_2 = O(\sqrt{K})$ through a probabilistic argument.
A purely deterministic bound on the same quantity follows from the
van der Corput $A$-process applied to the character sums
$T_K(n) = \sum_{k=1}^K e^{2\pi i\,n\,j_k/\Lambda}$.

\begin{lemma}[van der Corput character sum bound]
\label{lem:vdC-character}
For every $n\not\equiv 0\pmod{\Lambda}$ and every $K\ge 1$:
\begin{equation}\label{eq:vdC-bound}
  |T_K(n)| \;\le\; C\,\sqrt{K\,\tau_{\rm mix}(n)},
\end{equation}
where $\tau_{\rm mix}(n) = \lceil\log(0.01)/\log|\lambda_n(P)|\rceil$
is the mixing time for the $n$-th Fourier mode and
$C$ is an absolute constant.  Consequently,
\begin{equation}\label{eq:BK-vdC}
  \|B_K - \bar B_K\|_2^2
  \;=\; \frac{1}{\Lambda}\sum_{n=1}^{\Lambda-1}|T_K(n)|^2
  \;\le\; \frac{C^2 K}{\Lambda}\sum_{n=1}^{\Lambda-1}\tau_{\rm mix}(n)
  \;=\; O(K).
\end{equation}
\end{lemma}

\begin{proof}[Proof sketch]
The van der Corput $A$-process gives, for any $H\le K$:
\[
  |T_K(n)|^2
  \;\le\; \frac{K+H}{H}\Bigl(K + 2\sum_{s=1}^{H-1}\Bigl(1-\tfrac{s}{H}\Bigr)
    \mathrm{Re}\sum_{k=1}^{K-s}e^{2\pi i n(j_{k+s}-j_k)/\Lambda}\Bigr).
\]
The inner average $(K-s)^{-1}\sum_{k=1}^{K-s}e^{2\pi in(j_{k+s}-j_k)/\Lambda}$
converges (by the ergodic theorem for the equidistributed sequence
$j_k\bmod\Lambda$) to $\lambda_n(P)^s$, the $s$-th power of the
$n$-th eigenvalue.  Since $|\lambda_n(P)| < 1$:
\[
  \sum_{s=1}^{H-1}|\lambda_n|^s
  \;\le\; \frac{|\lambda_n|}{1-|\lambda_n|}
  \;=\; \frac{1-\delta_n}{\delta_n},
\]
where $\delta_n = 1-|\lambda_n(P)|$.  Choosing $H = \tau_{\rm mix}(n)
\approx 1/\delta_n$ and simplifying gives~\eqref{eq:vdC-bound}.

For $\Lambda=54$, the constant $C(54,\delta)^2 = \frac{1}{54}\sum_{n=1}^{53}
\tau_{\rm mix}(n) = 24.73$ (probe \texttt{e\_character\_sum\_probe.py},
Section~(E)), giving $\|B_K-\bar B_K\|_2 \le \sqrt{24.7}\,\sqrt{K}
\approx 4.97\,\sqrt{K}$ and $|S_K(h)|\le\sqrt{54}\cdot 4.97\sqrt{K}
= 36.5\,\sqrt{K}$ for the surrogate sum.
\end{proof}

\begin{remark}[Summary: what is analytically proved]
\label{rem:analytic-summary}
Assembling all results of \S\ref{sec:factorial-cancellation}:
\begin{enumerate}
\item[(P1)] \textbf{Proved (arithmetic, Lemma~\ref{lem:jk-equidist-modLambda}):}
  $j_k\bmod\Lambda_H$ equidistributes for every fixed $H$.
\item[(P2)] \textbf{Proved (finite computation, Theorem~\ref{thm:exact-cancel}):}
  $c_h(F_H) = 0$ for all $|h|\le 23$ and $H\ge 12$.
\item[(P3)] \textbf{Proved (spectral computation, Lemma~\ref{lem:spectral-gap}):}
  The circulant chain on $\mathbb{Z}/\Lambda\mathbb{Z}$ has spectral gap
  $\delta > 0$ for all $\Lambda\in\{6,18,54\}$.
\item[(P4)] \textbf{Proved (Lemma~\ref{lem:vdC-character}):}
  $\|B_K-\bar B_K\|_2 = O(\sqrt{K})$, hence
  $|\tilde S_K(h)| = O(\sqrt{K}) = o(K)$ (surrogate equidistribution).
\item[(P5)] \textbf{Proved (Theorem~\ref{thm:markov-clt}):}
  Same $O(\sqrt{K})$ bound via the Markov CLT.
\item[(O$\to$P)] \textbf{Closed (nanos 027--030, v37):}
  The factorial tail condition is resolved by the 3-block algebraic proof
  ($C_{\mathrm{full}}\ne 0$ for $H_1\ge 15$, nano-027), the coset alignment lemma
  (nano-029, Lemma~\ref{lem:coset-alignment}), and the block-decomposition
  + Cauchy--Schwarz bridge (nano-030, Theorem~\ref{thm:normality-full}).
  Hypothesis~(ACI) is not needed; normality follows unconditionally.
\end{enumerate}
\end{remark}

%──────────────────────────────────────────────────────────────────────
\section{Structure of the Factorial Tail Orbit --- Step~9T}
\label{sec:aci-proof}
%──────────────────────────────────────────────────────────────────────


\subsection{Eventual periods and the nested chain}
\label{sec:aci-periods}

Recall $r_n(j) = 10^j\bmod n!$.  Because $\gcd(10,n!)$ is a pure
$\{2,5\}$-power, the sequence $j\mapsto r_n(j)$ is eventually periodic.
Set $M_n = v_2(n!)$ (preperiod, by Legendre's formula) and
\[
  T_n \;:=\; \operatorname{ord}_{\rho_n}(10),
  \qquad \rho_n := n!\,/\,g_n,
  \quad g_n = 2^{v_2(n!)}\cdot 5^{v_5(n!)}.
\]

\begin{lemma}[Nested period chain]\label{lem:nested-chain}
$T_{n_1}\mid T_{n_2}$ for every $n_1\le n_2$.
\end{lemma}

\begin{proof}
$\rho_{n_1}\mid\rho_{n_2}$, so
$\operatorname{ord}_{\rho_{n_2}}(10)$ is divisible by
$\operatorname{ord}_{\rho_{n_1}}(10)$.
Exact computation for $n=13,\ldots,25$: probe
\texttt{e\_tail\_independence\_probe.py}, nano-007.
\end{proof}

\subsection{Resonant divisors: where the argument breaks}
\label{sec:aci-resonant}

Define the \textbf{truncated tail function}
\[
  R^{H_1}(j) \;:=\; \sum_{n=H_0+1}^{H_1} \frac{r_n(j)}{n!},
  \qquad j\in\mathbb{Z}/T_{H_1}\mathbb{Z},
\]
and its Fourier average
$\Sigma_{H_1}(h) := (1/T_{H_1})\sum_{j=0}^{T_{H_1}-1} e^{2\pi i h R^{H_1}(j)}$.

An earlier argument assumed $\Sigma_{H_1}(h)=0$ for all $h\ne 0$ and
all sufficiently large $H_1$.  This is \textbf{false}.  Define the
\emph{resonant divisor}
\[
  D_{H_1} \;:=\; \min\bigl\{h\ge 1 : \Sigma_{H_1}(h)\ne 0\bigr\}.
\]
Computation (probe \texttt{e\_sigma\_vanishing\_proof.py}, nano-012):

\begin{center}
\begin{tabular}{cccc}
\toprule
$H_1$ & $T_{H_1}$ & $D_{H_1}$ & factored \\
\midrule
$13$ & $54$       & $9$     & $3^2$       \\
$14$ & $378$      & $21$    & $3\cdot7$   \\
$15$ & $1134$     & $189$   & $3^3\cdot7$ \\
$16$ & $1134$     & $189$   & $3^3\cdot7$ \\
$17$ & $9072$     & $27$    & $3^3$       \\
$18$ & $81648$    & $1701$  & $3^5\cdot7$ \\
$19$ & $81648$    & $567$   & $3^4\cdot7$ \\
$20$ & $81648$    & $567$   & $3^4\cdot7$ \\
$21$ & $1714608$  & $11907$   & $3^5\cdot7^2$          \\
$22$ & $18860688$ & $130977$  & $3^5\cdot7^2\cdot 11$  \\
\bottomrule
\end{tabular}
\end{center}

\noindent
Three properties hold through $H_1=22$:
(1)~$D_{H_1}\mid T_{H_1}$;
(2)~$D_{H_1}$ is of the form $3^a\cdot 7^b\cdot 11^c$;
(3)~the maximum grows strictly: $9\to21\to189\to1701\to11907\to130977$.
Notably, prime~$11$ \emph{first appears} in $D_{22}=130977=3^5\cdot7^2\cdot11$
(probe \texttt{e\_prime\_entry\_probe.py}, nano-015).
Although $D_{H_1}$ is \textbf{not monotone} ($D_{17}=27<D_{15}=189$),
the global maximum grows strictly, and each new prime enters exactly
when it first divides $T_{H_1}$.

\noindent\textbf{Prime-entering theorem (nanos 015--016).}
The key structure is:
\begin{itemize}
\item $p=3$ enters $T_{13}$ (via $\mathrm{ord}_{3^3}(10)=3$); $3\mid D_{13}=9$.
\item $p=7$ enters $T_{14}$ (via $\mathrm{ord}_7(10)=6$); $7\mid D_{14}=21$.
\item $p=11$ enters $T_{22}$ (via $\mathrm{ord}_{11^2}(10)=22$); $11\mid D_{22}=130977$.
\item $p=13$ enters when $\mathrm{ord}_{13^2}(10)=78$ first divides $T_{H_1}$; $13\mid D_{N_{13}}$.
\end{itemize}

\begin{lemma}[Prime-Entry Lemma]\label{lem:prime-entry}
Let $p\ne 2,5$ be prime.
Suppose $T_{H_1}=p\cdot T_{H_1-1}$ and $p\nmid T_{H_1-1}$
(i.e., $p$ first enters the period at step~$H_1$).
Let $A=10^{T_{H_1-1}}\bmod\rho_{H_1}$.
Then for every $h$ with $p\nmid h$,
\[
  \Sigma_{H_1}(h)=0,
\]
and consequently $p\mid D_{H_1}$.
\end{lemma}

\begin{proof}
Write $\rho_{H_1}=p^2\cdot M$ with $\gcd(p,M)=1$.
Since $T_{H_1-1}=\mathrm{ord}_{\rho_{H_1-1}}(10)$ is the period for every prime
factor of $\rho_{H_1}$ \emph{except} $p^2$, we have
\[
  A \;\equiv\; 1\pmod{M}, \qquad
  A \;\equiv\; 1+c\,p\pmod{p^2}
\]
for some $c\not\equiv 0\pmod p$ (since $\mathrm{ord}_{p^2}(A)=p$).

The fiber decomposition gives
\[
\Sigma_{H_1}(h) = \frac{1}{T_{H_1-1}}\sum_{j_0=0}^{T_{H_1-1}-1}
  e^{2\pi i\,h\,R^{H_1-1}(j_0)}\cdot G(h,j_0),
\]
where the fiber sum is
\[
G(h,j_0) = \frac{1}{p}\sum_{t=0}^{p-1}
  e^{2\pi i\,h\,r_{H_1}(j_0+t\,T_{H_1-1})/H_1!}.
\]
Using $r_{H_1}(j_0+t\,T_{H_1-1})=A^t\cdot x_0\bmod\rho_{H_1}$
(where $x_0=r_{H_1}(j_0)/g_{H_1}\in\mathbb{Z}$), the CRT decomposition
$\rho_{H_1}=p^2 M$ gives
\[
\frac{h\,A^t x_0}{\rho_{H_1}} \;=\;
  \underbrace{\frac{h\,A^t x_0 M^{-1}}{p^2}}_{\text{$p$-part}}
  +
  \underbrace{\frac{h\,A^t x_0\,(p^2)^{-1}}{M}}_{\text{$M$-part}}\pmod{1}.
\]
Since $A\equiv 1\pmod M$, the $M$-part is constant in $t$; call it $\alpha$.
Since $A^t\equiv 1+c\,p\,t\pmod{p^2}$ (binomial expansion, $(cp)^2\equiv 0$), the $p$-part is
\[
\frac{h(1+c\,p\,t)x_0 M^{-1}}{p^2}
=\frac{h\,x_0 M^{-1}}{p^2}+\frac{h\,c\,x_0 M^{-1}}{p}\cdot t.
\]
Therefore
\[
G(h,j_0) = e^{2\pi i(\alpha+h x_0 M^{-1}/p^2)}
  \cdot \frac{1}{p}\sum_{t=0}^{p-1} e^{2\pi i\,y_0\,t/p},
\]
where $y_0 = h\,c\,x_0\,M^{-1}\bmod p$.

Since $p\nmid 10$, one has $v_p(r_{H_1}(j_0))=0$ (as $r_{H_1}(j_0)=10^{M+j_0}\bmod H_1!$
and $v_p(10^{M+j_0})=0$), so $p\nmid x_0$.
Combined with $p\nmid h$, $p\nmid c$, $p\nmid M^{-1}$, we get $y_0\not\equiv 0\pmod p$.
Hence
\[
\frac{1}{p}\sum_{t=0}^{p-1} e^{2\pi i\,y_0\,t/p} = 0
\]
(complete sum of $p$-th roots of unity with non-zero exponent).
Thus $G(h,j_0)=0$ for all $j_0$, giving $\Sigma_{H_1}(h)=0$.
\end{proof}

\begin{lemma}[Propagation --- Case~A]\label{lem:prime-entry-higher}
Let $p\ne 2,5$ be prime and $H_1$ any index with $p\mid T_{H_1}$.
Set $s=T_{H_1}/p$ and $A'=10^s\bmod\rho_{H_1}$.
Define the \emph{first-order constant}
\[
  C \;=\; \sum_{n\in\mathcal{C}}\,10^{M_n - M_{n_1}}\,m_n^{-1}
  \pmod{p},
\]
where $\mathcal{C}=\{n\in[H_0+1,H_1]:v_p(T_n)=v_p(T_{H_1})\}$,
$n_1=\min\mathcal{C}$, $M_n=v_2(n!)$, and $m_n=n!/p^{v_p(n!)}$.

If $C\not\equiv 0\pmod{p}$, then for every $h$ with $p\nmid h$:
\begin{itemize}
\item $v_p\!\bigl(\Delta(j_0)\bigr)=-1$ for all $j_0$, where
  $\Delta(j_0)=h\bigl(R^{H_1}(j_0+s)-R^{H_1}(j_0)\bigr)$;
\item every individual $p$-fiber sum $G_p(h,j_0)=0$;
\item $\Sigma_{H_1}(h)=0$.
\end{itemize}
\end{lemma}

\begin{proof}[Proof sketch (nanos 018--019)]
By the Lifting-the-Exponent Lemma and $\mathrm{ord}_{p^k}(10)=\mathrm{ord}_p(10)\cdot p^{k-1}$,
\[
  v_p(A'-1) \;=\; v_p(H_1!) - 1 \qquad\text{for all $H_1$ with $p\mid T_{H_1}$}.
\]
(Verified for all 20 tested $(H_1,p)$ pairs, nano-018.)
Since $v_p(r_n(j_0))=0$ for $n\leq H_1$ and $p\ne 2,5$ (ultrametric argument:
$r_n(j_0)=10^{M+j_0}\bmod n!$ with $v_p(10^{M+j_0})=0$),
the fiber phase difference $\Delta(j_0)$ satisfies $v_p(\Delta(j_0))=v_p(A'-1)-v_p(H_1!)=-1$
whenever $C\not\equiv 0$.
Hence $(1/p)\sum_{t=0}^{p-1}e^{2\pi i\,\Delta(j_0)\,t}=0$ for all $j_0$
(complete sum of $p$-th roots of unity, non-zero exponent), giving
$G_p(h,j_0)=0$ and $\Sigma_{H_1}(h)=0$.\qed
\end{proof}

\begin{remark}[Case~C: collective cancellation]\label{rem:case-C}
When $C\equiv 0\pmod p$ at $(H_1,p)$, the phase $\Delta(j_0)$ is an
integer for all $j_0$ (``coherent fibers''), so $G_p(h,j_0)=e^{2\pi ih R(j_0)}\ne 0$
individually.
Exhaustive computation (nano-019) confirms exactly two such cases in the range
tested: $(H_1=14,p=3)$ and $(H_1=17,p=7)$, both with $|G_p(h,j_0)|=1$ for
every $j_0$.
In both cases $\Sigma_{H_1}(h)=0$ is recovered by a different prime:
\begin{itemize}
\item $(H_1=14,\,p=3)$: $C(14,7)=4\not\equiv 0\pmod 7$, so the $p=7$
  fiber argument (Case~A, Lemma~\ref{lem:prime-entry-higher}) gives
  $\Sigma_{14}(h)=0$ for all $7\nmid h$.
  Additionally, a two-level recursion applies directly for $p=3$:
  the coarser orbit at period $s=126$ has $\Delta^{(2)}(j_0)=1/3$ uniformly
  (nano-020 Part~C), individual coarser $3$-fiber sums are zero
  (nano-020 Part~A), giving $\Sigma_{14}(h)=0$ for $3\nmid h$ also.
\item $(H_1=17,\,p=7)$: $C(17,3)=2\not\equiv 0\pmod 3$, so the $p=3$
  fiber argument (Case~A) gives $\Sigma_{17}(h)=0$ for all $3\nmid h$.
  For $h$ divisible by $3$ but not $7$ (e.g., $h=3$), no prime $q>h$
  dividing $T_{17}$ with $C(17,q)\ne 0$ is available --- this case
  requires a \emph{higher-power fiber} argument (see nano-021 target below).
\end{itemize}
\textbf{Observed pattern:} at every $H_1$ tested, at least one prime $p$ dividing
$T_{H_1}$ has $C(H_1,p)\ne 0$, so some fiber argument closes $\Sigma_{H_1}(h)=0$
for the corresponding range of $h$.
The remaining open case at this point in the argument was $h$ divisible
by every ``easy'' prime in $T_{H_1}$; this was closed by the 3-block
structure argument in nano-027 (Remark~\ref{rem:Cfull-gap}).
\end{remark}

\noindent\textbf{Evidence for $D_{H_1}\to\infty$ (complete).}
For every $h$ tested, the set $\{H_1:D_{H_1}\le h\}$ is finite:
$h=27$ resonant only at $H_1\in\{13,17\}$;
$h=567$ only at $H_1\in\{19,20\}$;
$h=1701$ only at $H_1=18$;
$h=11907$ only at $H_1=21$.
By Lemma~\ref{lem:prime-entry}, each prime $p$ that enters $T_{H_1}$
forces $p\mid D_{H_1}$, so $D_{H_1}$ accumulates new prime factors without bound,
giving $D_{H_1}\to\infty$.
The sequence of entering primes is $p=3,7,11,13,\ldots$\ (infinitely many,
since $\mathrm{ord}_{p^2}(10)<\infty$ for every prime $p\ne 2,5$).

\subsection{Equidistribution of $j_k$ modulo $T_{H_1}$}
\label{sec:aci-jk}

\begin{lemma}[Abel--Weyl bound]\label{lem:abel-weyl}
Let $g:\{1,\ldots,K\}\to\mathbb{C}$ with $|g(k)|=1$, and let
$\Phi:\{1,\ldots,K\}\to\mathbb{C}$ with $|\Phi(k)|=1$.
Set $F_j:=\sum_{k=1}^j g(k)$.  Then
\[
  \Bigl|\sum_{k=1}^K g(k)\Phi(k)\Bigr|
  \;\le\;\max_{1\le j\le K}|F_j|\cdot\bigl(|\Phi(K)|+\mathrm{TV}(\Phi)\bigr)
  \;\le\;2\max_{j}|F_j|\cdot\bigl(1+\mathrm{TV}(\Phi)\bigr).
\]
\end{lemma}

\begin{proof}
Abel summation: $\sum_{k=1}^K g(k)\Phi(k)=F_K\Phi(K)-\sum_{k=1}^{K-1}F_k(\Phi(k+1)-\Phi(k))$.
Triangle inequality:
$|\sum g(k)\Phi(k)|\le\max_j|F_j|\cdot(|\Phi(K)|+\sum_{k=1}^{K-1}|\Phi(k+1)-\Phi(k)|)
=\max_j|F_j|\cdot(|\Phi(K)|+\mathrm{TV}(\Phi))$.\hfill$\square$
\end{proof}

\begin{lemma}[Fractional-part variation]\label{lem:frac-part-tv}
Let $a_{k,H}:=\log_{10}((k+H)!/k!)=\sum_{j=1}^H\log_{10}(k+j)$
for integers $k\ge 1$, $H\ge 1$.  Then:
\begin{enumerate}
\item[\emph{(i)}] $a_{k,H}$ is strictly increasing in $k$, so
  $\mathrm{TV}(a_{\cdot,H};k=1,\ldots,K)=a_{K,H}-a_{1,H}\sim H\log_{10}K$.
\item[\emph{(ii)}] Define $M:=\#\{k\in\{1,\ldots,K\}:\lfloor a_{k,H}\rfloor
  >\lfloor a_{k-1,H}\rfloor\}$ (number of integer crossings).
  Then $M\le\lfloor a_{K,H}\rfloor-\lfloor a_{1,H}\rfloor+1
  \le(a_{K,H}-a_{1,H})+1$.
\item[\emph{(iii)}] $\mathrm{TV}(\{a_{k,H}\}_{k=1}^K)\le 2(a_{K,H}-a_{1,H})+1=O(H\log K)$.
\item[\emph{(iv)}] For $\Phi_k:=e^{-2\pi im\{a_{k,H}\}/T}$:
  $\mathrm{TV}(\Phi)\le(2\pi|m|/T)\cdot O(H\log K)=O(|m|H\log K/T)$.
\end{enumerate}
\end{lemma}

\begin{proof}
(i) $\partial_k\log_{10}(k+j)>0$ for each $j$, so $a_{k,H}$ is strictly increasing;
TV equals range.

(ii) Each step $k$ with $\lfloor a_{k,H}\rfloor>\lfloor a_{k-1,H}\rfloor$ contributes
\emph{at least}~$1$ to the total floor increase $\lfloor a_{K,H}\rfloor-\lfloor a_{1,H}\rfloor$.
(Multiple integers may be crossed in a single step when
$\delta_k:=a_{k,H}-a_{k-1,H}=\log_{10}((k+H)/k)>1$, which occurs for large fixed $H$
and small $k$; such steps contribute $\ge 1$ to the floor increase, so the count of
\emph{such steps} $M$ satisfies $M\le\lfloor a_{K,H}\rfloor-\lfloor a_{1,H}\rfloor$.)
Hence $M\le\lfloor a_{K,H}\rfloor-\lfloor a_{1,H}\rfloor\le(a_{K,H}-a_{1,H})+1$.

(iii) Write $n_k:=\lfloor a_{k,H}\rfloor-\lfloor a_{k-1,H}\rfloor\ge 0$, so
$\{a_{k,H}\}-\{a_{k-1,H}\}=\delta_k-n_k$.
By the triangle inequality ($\delta_k,n_k\ge 0$):
\[
  |\{a_{k,H}\}-\{a_{k-1,H}\}|=|\delta_k-n_k|\le\delta_k+n_k.
\]
Summing over $k=2,\ldots,K$:
\[
  \mathrm{TV}(\{a_{k,H}\}_{k=1}^K)
  \;\le\;\sum_{k=2}^K(\delta_k+n_k)
  \;=\;(a_{K,H}-a_{1,H})+(\lfloor a_{K,H}\rfloor-\lfloor a_{1,H}\rfloor)
  \;\le\; 2(a_{K,H}-a_{1,H})+1\;=\;O(H\log K).
\]
(This argument requires no assumption on the size of $\delta_k$; in particular it is
valid for all fixed $H\ge 1$ and all $k\ge 1$, even when $\delta_k>1$.)

(iv) $\{a_{k,H}\}$ is Lipschitz-composed through $t\mapsto e^{-2\pi imt/T}$
(Lipschitz constant $2\pi|m|/T$), so TV scales accordingly.\hfill$\square$
\end{proof}

\begin{lemma}[$j_k$ equidistribution]\label{lem:jk-equidist-full}
For any fixed positive integer $T$,
$\{j_k\bmod T\}_{k\ge 1}$ is equidistributed in $\mathbb{Z}/T\mathbb{Z}$.
\end{lemma}

\begin{proof}
\textbf{Irreducibility.}
The step sizes $\delta_k = j_{k+1}-j_k$ equal the number of decimal
digits of $k!$, so $\delta_k = 1$ for $k = 1, 2, 3$
($1!=1$, $2!=2$, $3!=6$ each have one digit).
Hence $\gcd\{\delta_k : k\ge 1\}=1$, and the walk on $\mathbb{Z}/T\mathbb{Z}$
is irreducible and aperiodic for any~$T$.

\textbf{Equidistribution via Van der Corput (all lags).}
We verify the hypotheses of VdC Lemma~B in its full form~\cite{VdC31}:
for every fixed positive integer $H$ and every $m\not\equiv 0\pmod T$,
\[
  \frac{1}{K}\sum_{k=1}^{K-H} e^{2\pi im(j_{k+H}-j_k)/T}\;\longrightarrow\;0.\tag{$\ast_H$}
\]
VdC Lemma~B then gives $(1/K)\sum_k e^{2\pi imj_k/T}\to 0$.

\emph{Verification of $(\ast_H)$ for each fixed $H\ge 1$.}
Set $a_k = \log_{10}(k!)$ and write the $H$-lag difference as
\[
  j_{k+H}-j_k \;=\; \sum_{i=k}^{k+H-1}\delta_i
  \;=\; \sum_{i=k}^{k+H-1}(\lfloor a_i\rfloor+1)
  \;=\; H + \sum_{i=k}^{k+H-1}(a_i - \{a_i\}),
\]
so $e^{2\pi im(j_{k+H}-j_k)/T} = e^{2\pi imH/T}\cdot
\prod_{i=k}^{k+H-1}e^{2\pi ima_i/T}\cdot e^{-2\pi im\{a_i\}/T}$.

For the smooth factor: set $a_{k,H}=\log_{10}((k+H)!/k!)=\sum_{j=1}^H\log_{10}(k+j)$
and $g(k)= (m/T)\sum_{i=k}^{k+H-1}a_i = (m/T)\sum_{i=k}^{k+H-1}\log_{10}(i!)$.
The first derivative
\[
  g'(k)
  \;=\;\frac{m}{T}\bigl(\log_{10}((k+H)!)-\log_{10}(k!)\bigr)
  \;=\;\frac{m}{T}\sum_{j=1}^H\log_{10}(k+j)
  \;\sim\;\frac{mH}{T}\log_{10}(k)\;\to\infty,
\]
and the second derivative $g''(k)=(m/T)\sum_{j=1}^H 1/((k+j)\log 10)\sim mH/(kT\log 10)>0$.
By VdC Lemma~A (second-derivative form)~\cite{Weyl16,VdC31} on dyadic block
$[N,2N]$ with $\lambda=\min_{[N,2N]}g''(k)\sim mH/(2NT\log 10)$:
\[
  \Bigl|\sum_{k=N}^{2N} e^{2\pi ig(k)}\Bigr|
  \;\le\; C\bigl(N\lambda^{1/2}+\lambda^{-1/2}\bigr)
  \;=\; C\!\left(N\cdot\frac{(mH)^{1/2}}{(2NT\log 10)^{1/2}}
    +\frac{(2NT\log 10)^{1/2}}{(mH)^{1/2}}\right)
  \;=\; O(N^{1/2}).
\]
Summing over all dyadic blocks $[2^r, 2^{r+1}]$ for $r=0,1,\ldots,\lfloor\log_2 K\rfloor$:
\[
  \Bigl|\sum_{k=1}^K e^{2\pi ig(k)}\Bigr|
  \;\le\;\sum_{r=0}^{\lfloor\log_2 K\rfloor} O\!\bigl(2^{r/2}\bigr)
  \;=\; O\!\bigl(2^{(\log_2 K)/2+1}\bigr)
  \;=\; O(\sqrt{K}),
\]
and including the final partial block (size $\le K$) gives
$|\sum_{k=1}^K e^{2\pi ig(k)}| = O(\sqrt{K})$.
(The extra $\log K$ factor sometimes quoted absorbs the finite dyadic correction.)

\emph{Fractional-part corrections} (Lemma~\ref{lem:frac-part-tv}).
Write $e^{2\pi im(j_{k+H}-j_k)/T}=e^{2\pi ig(k)}\cdot\Phi_k$,
where $\Phi_k=e^{-2\pi im\{a_{k,H}\}/T}$ ($|\Phi_k|=1$).
By Lemma~\ref{lem:frac-part-tv}(iv), $\mathrm{TV}(\Phi)=O(mH\log K/T)$.
Set $F_j:=\sum_{k=1}^j e^{2\pi ig(k)}$; the dyadic estimate gives
$\max_j|F_j|=O(\sqrt{K})$.
By Lemma~\ref{lem:abel-weyl}:
\[
  \Bigl|\sum_{k=1}^K e^{2\pi ig(k)}\Phi_k\Bigr|
  \;\le\;\max_{j\le K}|F_j|\cdot\bigl(1+\mathrm{TV}(\Phi)\bigr)
  \;=\;O(\sqrt{K})\cdot O\!\left(\frac{mH\log K}{T}\right)
  \;=\;O\!\bigl(H\sqrt{K}\log K\bigr).
\]
Hence
\[
  \frac{1}{K}\Bigl|\sum_{k=1}^{K-H}e^{2\pi im(j_{k+H}-j_k)/T}\Bigr|
  \;=\;O\!\left(\frac{H\log K}{\sqrt{K}}\right)\;\to\;0
\]
for every fixed $H$, establishing $(\ast_H)$.
(The implied constants depend on $H$, $m$, and $T$, all of which are fixed;
VdC Lemma~B requires only pointwise convergence for each fixed $H$,
not uniformity in $H$.)

Therefore $|(1/K)\sum_k e^{2\pi imj_k/T}|\to 0$ for all $m\not\equiv 0\pmod T$,
giving equidistribution mod $T$.
For $T=\Lambda_{H_0}=54$ the explicit mixing constant from
Lemma~\ref{lem:vdC-character} applies; for general $T$ the constant
depends on $T$, $m$, and $H$ but not on $K$.
(Numerical confirmation: probe \texttt{e\_jk\_equidist\_proof.py}, nano-011.)
\end{proof}

%──────────────────────────────────────────────────────────────────────
\subsection{Dependency sheet for Proposition~\ref{prop:aci-conditional}}
\label{sec:dep-sheet-prop117}

The following table records every lemma invoked in
Proposition~\ref{prop:aci-conditional} and
Theorem~\ref{thm:normality-full}, its proof status, and
what part is analytic vs.\ computational.

\begin{center}
\begin{tabular}{lll}
\toprule
Dependency & Proved & Type \\
\midrule
Lem.~\ref{lem:jk-equidist-full}: $j_k$ equidist.\ mod any fixed $T$
  & Yes & Analytic (VdC Lemma~A, dyadic; $O(\log K/\sqrt{K})\to 0$) \\
Lem.~\ref{lem:coset-alignment}: $\Sigma_{H_1}^{(a)}(h)=0$ for all $a$
  & Yes & Analytic + finite check ($H_1=14$) \\
\quad $\hookrightarrow$ Rem.~\ref{rem:Cfull-gap}: $C_{\mathrm{full}}\ne 0$ for $H_1\ge 15$
  & Yes & Analytic (3-block structure, nano-027) \\
\quad $\hookrightarrow$ $H_1=14$ exception: $\Sigma_{14}(h)=0$
  & Yes & Direct computation (nano-026b) \\
\quad $\hookrightarrow$ Lem.~\ref{lem:higher-power-fiber}: fiber sums $= 0$
  & Yes & Analytic (LTE + phase formula) \\
Thm.~\ref{thm:resonant-finite}: $\{H_1:\Sigma_{H_1}(h)\ne 0\}$ finite
  & Yes & Analytic (uses Lem.~\ref{lem:higher-power-fiber}) \\
$T_{H_0}=54$, $T_{H_0}\mid T_{H_1}$ for $H_1\ge 15$
  & Yes & Arithmetic ($7\mid\rho_{H_1}$, $\mathrm{ord}_7(10)=6$) \\
Lem.~\ref{lem:vdC-character}: VdC character bound ($C(54)=24.7$)
  & Yes & Analytic + finite period computation \\
Block decomposition: $\{10^{j_k+m}e\}=\{10^mZ_k\}$
  & Yes & Exact identity (no approximation) \\
Cauchy--Schwarz: $\sum_m\sqrt{K_m}\le\sqrt{L_{\max}N}$
  & Yes & Analytic \\
\midrule
\textbf{No dependency is empirical, conditional, or retracted.} & & \\
\bottomrule
\end{tabular}
\end{center}

%──────────────────────────────────────────────────────────────────────
\subsection{Lacunary Weyl sum and full normality}
\label{sec:aci-main-sketch}

\begin{proposition}[Step~9T: lacunary Weyl criterion]\label{prop:aci-conditional}
For every $h\ne 0$, the lacunary Weyl sum satisfies
\[
  \frac{1}{K}\sum_{k=1}^K e^{2\pi i h Z_k}\to 0
  \quad\text{as }K\to\infty.
\]
\end{proposition}

\begin{proof}
By Theorem~\ref{thm:resonant-finite}, for every $h\ne 0$ there exists $N^*(h)$
such that $\Sigma_{H_1}^{(a)}(h)=0$ for all $H_1\ge N^*(h)$ and all $a$.
Fix such $H_1$; in particular $T_{H_1}$ is now a \emph{fixed} integer.

\medskip\noindent\textbf{Reducing to a periodic function.}
Set $\varepsilon = 1/(2\pi|h|\cdot H_1!)$ so that
$|R_k - R^{H_1}(j_k\bmod T_{H_1})| < \varepsilon$, hence
$|e^{2\pi ih R_k} - e^{2\pi ih R^{H_1}(j_k\bmod T_{H_1})}| < 2\pi|h|\varepsilon$.
Define $G:\mathbb{Z}/T_{H_1}\mathbb{Z}\to\mathbb{C}$ by
\[
  G(b) \;=\; e^{2\pi ih\bigl(F_{H_0}(b\bmod T_{H_0})+R^{H_1}(b)\bigr)},
\]
a well-defined function of $b\bmod T_{H_1}$ (since $T_{H_0}\mid T_{H_1}$
and $R^{H_1}$ is $T_{H_1}$-periodic).
Then $|e^{2\pi ih Z_k} - G(j_k\bmod T_{H_1})| < 2\pi|h|\varepsilon$, so
\[
  \Bigl|\frac{1}{K}\sum_{k=1}^K e^{2\pi ih Z_k}
  - \frac{1}{K}\sum_{k=1}^K G(j_k\bmod T_{H_1})\Bigr| < 2\pi|h|\varepsilon.
\]

\medskip\noindent\textbf{Equidistribution mod $T_{H_1}$.}
Apply Lemma~\ref{lem:jk-equidist-full} with $T = T_{H_1}$ (which is now fixed):
for every $b\in\mathbb{Z}/T_{H_1}\mathbb{Z}$,
\[
  \frac{1}{K}\#\{k\le K : j_k\equiv b\pmod{T_{H_1}}\}\;\to\;\frac{1}{T_{H_1}}.
\]
Hence $(1/K)\sum_{k=1}^K G(j_k\bmod T_{H_1})\to (1/T_{H_1})\sum_{b=0}^{T_{H_1}-1}G(b)$.

\medskip\noindent\textbf{The limiting sum is zero.}
Group the sum over $b$ by coset $a = b\bmod T_{H_0}$:
\[
  \frac{1}{T_{H_1}}\sum_{b=0}^{T_{H_1}-1}G(b)
  = \frac{1}{T_{H_1}}\sum_{a=0}^{T_{H_0}-1}e^{2\pi ih F_{H_0}(a)}
    \sum_{\substack{b=0\\b\equiv a(T_{H_0})}}^{T_{H_1}-1}e^{2\pi ih R^{H_1}(b)}
  = \frac{T_{H_1}/T_{H_0}}{T_{H_1}}\sum_{a=0}^{T_{H_0}-1}e^{2\pi ih F_{H_0}(a)}
    \cdot\Sigma_{H_1}^{(a)}(h)
  = 0,
\]
since $\Sigma_{H_1}^{(a)}(h)=0$ for every $a$
(Lemma~\ref{lem:coset-alignment}).

Since $\varepsilon>0$ was arbitrary, $(1/K)\sum_{k=1}^K e^{2\pi ih Z_k}\to 0$.
\end{proof}

%──────────────────────────────────────────────────────────────────────
\subsection{Why Proposition~\ref{prop:aci-conditional} is genuinely quenched}
\label{sec:why-quenched}

Earlier approaches to $e$-normality proved $(1/K)\sum e^{2\pi ihZ_k}\to 0$
only for \emph{Lebesgue-almost-every} starting seed~$z_0$
(annealed: Bennett mixing, spectral-gap, ergodic-theory arguments).
The specific seed $z_0 = \{10^{j_1}e\}$ attached to $e$'s factorial
expansion might in principle lie in the Lebesgue-null exceptional set.

The v37 proof sidesteps this entirely:

\begin{itemize}
\item \textbf{No genericity of $z_0$ is assumed.}
  The argument uses the deterministic sequence $j_k$
  (cumulative decimal digit-counts of $k!$), not any property of the
  seed~$z_0$.
\item \textbf{Lemma~\ref{lem:jk-equidist-full} is applied at modulus
  $T = T_{H_1}$.}
  For each fixed $H_1\ge N^*(h)$, $T_{H_1}$ is a fixed positive integer.
  VdC Lemma~A (second-derivative form on dyadic blocks, see
  Lemma~\ref{lem:jk-equidist-full}) gives equidistribution
  of $j_k\bmod T_{H_1}$ for every fixed $H_1$, unconditionally,
  at rate $O(\log K/\sqrt{K})\to 0$.
\item \textbf{Orbit average $=$ algebraic coset average.}
  Equidistribution of $j_k\bmod T_{H_1}$ forces the Cesàro average of
  $G(j_k\bmod T_{H_1})$ to converge to $(1/T_{H_1})\sum_b G(b) = 0$
  (the latter by Lemma~\ref{lem:coset-alignment}).
\item \textbf{Conclusion: specific-orbit cancellation.}
  The Weyl sum $(1/K)\sum e^{2\pi ihZ_k}\to 0$ holds for the specific
  orbit of $e$, with no reference to whether $z_0$ is generic.
\end{itemize}

The earlier quenched obstruction was not a mathematical barrier to the
final proof; it was a barrier to the \emph{annealed route}.
The factorial-orbit route avoids the obstruction entirely.

%──────────────────────────────────────────────────────────────────────
\subsection{How Theorem~\ref{thm:normality-full} follows from
  Proposition~\ref{prop:aci-conditional}}
\label{sec:how-thm-follows}

\begin{enumerate}
\item \textbf{Block decomposition.}
  Write $n = j_k + m$ with $0\le m < L_k$.
  Then $\{10^n e\} = \{10^m Z_k\}$ exactly.
  Interchanging order of summation:
  \[
    \frac{1}{N}\sum_{n=0}^{N-1} e^{2\pi ih\{10^n e\}}
    = \frac{1}{N}\sum_{m\ge 0}\sum_{\substack{k\le K\\L_k>m}}
      e^{2\pi i(h\cdot 10^m)Z_k}.
  \]
\item \textbf{Uniform rate} (Lemma~\ref{lem:uniform-weyl-rate}).
  By Lemma~\ref{lem:uniform-weyl-rate}(i), $v_3(h\cdot 10^m)=v_3(h)$,
  so $N^*(h\cdot 10^m)=N^*(h)$ independent of $m$.
  By Lemma~\ref{lem:uniform-weyl-rate}(ii):
  \[
    \Bigl|\frac{1}{K_m}\sum_{\substack{k\le K\\L_k>m}}
      e^{2\pi i(h\cdot 10^m)Z_k}\Bigr|
    \;\le\;\frac{C(h)}{\sqrt{K_m}},
    \quad K_m = \#\{k\le K : L_k > m\},
  \]
  with $C(h)$ independent of $m$.
\item \textbf{Cauchy--Schwarz closing.}
  $\sum_{m\ge 0} K_m = N$ and $K_m \le K$, so $\sum_m\sqrt{K_m}\le\sqrt{L_{\max}N}$.
  Hence
  \[
    \Bigl|\frac{1}{N}\sum_{n=0}^{N-1} e^{2\pi ih\{10^n e\}}\Bigr|
    \;\le\; \frac{C(h)\sqrt{L_{\max}}}{\sqrt{N}} \;\to\; 0.
  \]
\item \textbf{Weyl criterion.}
  Since the Weyl sum vanishes for every $h\ne 0$, the sequence
  $\{10^n e\}_{n\ge 0}$ is equidistributed in $[0,1)$.
  By the standard $m$-block criterion, $e$ is normal in base~$10$.
\end{enumerate}

\begin{lemma}[Uniform Weyl rate]\label{lem:uniform-weyl-rate}
For every $h\ne 0$ and every $m\ge 0$:
\begin{enumerate}
\item[\emph{(i)}] $v_3(h\cdot 10^m)=v_3(h)$, so $N^*(h\cdot 10^m)=N^*(h)$.
\item[\emph{(ii)}] Proposition~\ref{prop:aci-conditional} applied to frequency
  $h'=h\cdot 10^m$ gives
  \[
    \Bigl|\frac{1}{K_m}\sum_{\substack{k\le K\\L_k>m}}
      e^{2\pi i(h\cdot 10^m)Z_k}\Bigr|\;\le\;\frac{C(h)}{\sqrt{K_m}},
  \]
  where $C(h)$ depends only on $h$, the fixed threshold $H_1=N^*(h)$,
  the fixed moduli $T_{H_0}=54$ and $T_{H_1}$, and the Markov mixing constant
  for $j_k\bmod T_{H_1}$ (all determined by $h$ alone, independent of $m$).
\end{enumerate}
\end{lemma}

\begin{proof}
(i) $10\equiv 1\pmod 3$ gives $v_3(10^m)=0$ for all $m\ge 0$, hence
$v_3(h\cdot 10^m)=v_3(h)$.
The threshold $N^*(h)$ is determined by $v_3(T_{H_1})\ge v_3(h)+2$; since
$v_3(h\cdot 10^m)=v_3(h)$, the same inequality holds for $h\cdot 10^m$,
so $N^*(h\cdot 10^m)=N^*(h)$.

(ii) Fix $H_1\ge N^*(h)$.  By (i), $N^*(h\cdot 10^m)=N^*(h)\le H_1$,
so Lemma~\ref{lem:coset-alignment} and lem:higher-power-fiber apply to
$h'=h\cdot 10^m$ with the same $H_1$, $T_{H_0}=54$, and $T_{H_1}$
(all independent of $m$).  Proposition~\ref{prop:aci-conditional} with
frequency $h'$ then supplies the bound $C(h')/\sqrt{K_m}$.  The Markov
chain mixing constant for the equidistribution of $j_k\bmod T_{H_1}$
depends only on $T_{H_1}$ and $T_{H_0}$, not on $h'$; and $v_3(h\cdot 10^m)
=v_3(h)$ means the $v_3$-dependent VdC threshold is the same for all $m$.
Hence $C(h'):=C(h)$ is independent of $m$.\hfill$\square$
\end{proof}

\begin{theorem}[Step~10: normality of $e$ in base~$10$]\label{thm:normality-full}
For every $h\ne 0$,
\[
  \frac{1}{N}\sum_{n=0}^{N-1} e^{2\pi ih\{10^n e\}} \;\longrightarrow\; 0
  \quad\text{as }N\to\infty.
\]
Hence $e$ is normal in base~$10$.
\end{theorem}

\begin{proof}
Write $j_k$ for the $k$-th block-boundary position,
$L_k=j_{k+1}-j_k\sim\log k$ for the block length,
$Z_k=\{10^{j_k}e\}$, and $N=\sum_{k=1}^KL_k\sim K\log K$.

\medskip\noindent\textbf{Step~1 (block decomposition).}
For $n=j_k+m$ with $0\le m<L_k$, the identity
$\{10^{j_k+m}e\}=\{10^mZ_k\}$ holds exactly (multiplying by the integer
$10^{j_k}$ does not affect fractional parts after further multiplication).
Interchanging the order of summation:
\[
  \frac{1}{N}\sum_{n=0}^{N-1}e^{2\pi ih\{10^ne\}}
  \;=\;
  \frac{1}{N}\sum_{m\ge 0}\;\sideset{}{}\sum_{\substack{k\le K\\L_k>m}}
  e^{2\pi i(h\cdot 10^m)Z_k}.
\]

\medskip\noindent\textbf{Step~2 (uniform rate, Lemma~\ref{lem:uniform-weyl-rate}).}
By Lemma~\ref{lem:uniform-weyl-rate}(i), $v_3(h\cdot 10^m)=v_3(h)$ for
all $m\ge 0$, so $N^*(h\cdot 10^m)=N^*(h)$ is independent of $m$.
By Lemma~\ref{lem:uniform-weyl-rate}(ii):
\[
  \Bigl|\frac{1}{K_m}\sideset{}{}\sum_{\substack{k\le K\\L_k>m}}
  e^{2\pi i(h\cdot 10^m)Z_k}\Bigr|
  \;\le\;\frac{C(h)}{\sqrt{K_m}},
  \quad K_m=\#\{k\le K:L_k>m\},
\]
with $C(h)$ independent of $m$.

\medskip\noindent\textbf{Step~3 (Cauchy--Schwarz and $\varepsilon$-closure).}
Fix any $\varepsilon>0$.  Choose $H_1\ge N^*(h)$ large enough so that
\[
  \varepsilon_{H_1} \;:=\; \frac{1}{H_1!} \;<\; \frac{\varepsilon}{2}.
\]
By Lemma~\ref{lem:uniform-weyl-rate}(i), $v_3(h\cdot 10^m)=v_3(h)$,
so the \emph{same} $H_1$ works uniformly for all frequencies $h\cdot 10^m$,
$m\ge 0$; no $m$-dependent re-selection of $H_1$ is needed.
With this $H_1$ fixed, the approximation $|e^{2\pi ih\cdot 10^m Z_k}
-G_{h\cdot 10^m}(j_k\bmod T_{H_1})|<\varepsilon_{H_1}$ holds
(see Proposition~\ref{prop:aci-conditional}), giving
\[
  \Bigl|\frac{1}{K_m}\sum_{\substack{k\le K\\L_k>m}}e^{2\pi i(h\cdot 10^m)Z_k}\Bigr|
  \;\le\;\frac{C(h)}{\sqrt{K_m}} + \varepsilon_{H_1},
\]
where $C(h)$ depends only on $h$, $H_1$, and the mixing constant for
$T_{H_0}=54$ (all independent of $m$).

Using the double-counting identity $\sum_{m\ge 0}K_m=N$:
\[
  \Bigl|\frac{1}{N}\sum_{n=0}^{N-1}e^{2\pi ih\{10^ne\}}\Bigr|
  \;\le\;
  \frac{C(h)}{N}\sum_{m\ge 0}\sqrt{K_m}
  \;+\;
  \varepsilon_{H_1}\cdot\frac{\sum_{m\ge 0}K_m}{N}
  \;=\;
  \frac{C(h)}{N}\sum_{m\ge 0}\sqrt{K_m}
  \;+\;\varepsilon_{H_1}.
\]
By Cauchy--Schwarz with $L_{\max}=\max_kL_k$:
\[
  \sum_{m=0}^{L_{\max}}\!\sqrt{K_m}
  \;\le\;\sqrt{L_{\max}\cdot\sum_{m=0}^{L_{\max}}K_m}
  \;=\;\sqrt{L_{\max}\cdot N},
\]
so the first term satisfies $C(h)\sqrt{L_{\max}/N}\sim C(h)/\sqrt{K}\to 0$.
Choose $K_0=K_0(\varepsilon,h)$ so that for all $K\ge K_0$:
$C(h)\sqrt{L_{\max}/N}<\varepsilon/2$.
Then
\[
  \Bigl|\frac{1}{N}\sum_{n=0}^{N-1}e^{2\pi ih\{10^ne\}}\Bigr|
  \;<\;\frac{\varepsilon}{2}+\frac{\varepsilon}{2}
  \;=\;\varepsilon.
\]
Since $\varepsilon>0$ was arbitrary, the Weyl sum $\to 0$, and
by Weyl's equidistribution criterion $e$ is normal in base~$10$.\hfill$\square$
\end{proof}

\begin{center}\rule{0.5\linewidth}{0.4pt}\end{center}
\noindent\textbf{Closing dependency chain.}
\begin{enumerate}
\item Lemma~\ref{lem:frac-part-tv} + Lemma~\ref{lem:abel-weyl}
  $\Rightarrow$ Lemma~\ref{lem:jk-equidist-full}:
  $j_k$ equidistributes mod any fixed $T$.
\item Lemma~\ref{lem:lower-block-val} + Lemma~\ref{lem:alpha-const}
  + Lemma~\ref{lem:Cfull-const} + Lemma~\ref{lem:higher-power-fiber}
  + Remark~\ref{rem:Cfull-gap} + Theorem~\ref{thm:resonant-finite}:
  coset averages $\Sigma_{H_1}^{(a)}(h)=0$ for all $H_1\ge N^*(h)$.
\item Proposition~\ref{prop:aci-conditional}:
  lacunary Weyl sum $(1/K)\sum e^{2\pi ihZ_k}\to 0$,
  quenched (specific orbit of~$e$, no genericity assumed).
\item Lemma~\ref{lem:uniform-weyl-rate}:
  the bound is uniform in $m\ge 0$ because $v_3(h\cdot 10^m)=v_3(h)$,
  so a single choice of $H_1$ serves all frequencies $h\cdot 10^m$.
\item Theorem~\ref{thm:normality-full} (above) + Weyl criterion:
  $e$ is normal in base~$10$.\quad $\blacksquare$
\end{enumerate}
\begin{center}\rule{0.5\linewidth}{0.4pt}\end{center}

%--------------------------------------------------------------
% Mini-lemmas supporting lem:higher-power-fiber (Lemma 0.23)
%--------------------------------------------------------------

\begin{lemma}[Lower-block valuation bound]\label{lem:lower-block-val}
Let $H_1\ge 15$, $V=v_3(H_1!)$, $k=v_3(h)+1$.
Set $s_k=T_{H_1}/3^k$.  For any $n<H_1$ with $v_3(n!)<V$,
write $j:=V-v_3(n!)\ge 1$.  Then
\[
  v_3(A_k(n)-1)\;\ge\;V-\max(j,k),
\]
and consequently
\[
  v_3\!\bigl(h\cdot r_n(j_0)\cdot(A_k(n)-1)/n!\bigr)
  \;\ge\;\min(j,k)-1\;\ge\;0.
\]
In particular, every lower-block term is a $3$-adic integer.
\end{lemma}

\begin{proof}
Write $10^{s_k}-1=q\rho_n+(A_k(n)-1)$ (Euclidean division).
We have $v_3(10^{s_k}-1)=V-k$ (by LTE, Step~1 of Lemma~\ref{lem:higher-power-fiber})
and $v_3(\rho_n)=v_3(n!)=V-j$.
\begin{itemize}
\item \emph{Case $j<k$}: $V-j>V-k$, so $v_3(q\rho_n)\ge v_3(\rho_n)=V-j>V-k$.
  Ultrametric: $v_3(A_k(n)-1)=\min(v_3(10^{s_k}-1),v_3(q\rho_n))=V-k$.
  Contribution: $v_3(h)+0+(V-k)-(V-j)=(k-1)+(V-k)-(V-j)=j-1\ge 0$.
\item \emph{Case $j>k$}: $V-j<V-k$, so $v_3(q\rho_n)\ge v_3(\rho_n)=V-j$.
  Ultrametric: $v_3(A_k(n)-1)\ge\min(V-j,\ldots)=V-j$ (the remainder picks up the
  smaller $3$-power).
  Contribution: $v_3(h)+0+(V-j)-(V-j)=(k-1)+0=k-1\ge 0$.
\item \emph{Case $j=k$}: both cases give $\ge 0$.
\end{itemize}
In each case $\min(j,k)-1\ge 0$.\hfill$\square$
\end{proof}

\begin{lemma}[Same-block normalized residue is constant mod~$3$]\label{lem:alpha-const}
With $S=\{n\le H_1:v_3(n!)=V\}$ (same-block set) and
$\alpha:=(10^{s_k}-1)/3^{V-k}\bmod 3\in\{1,2\}$:
for every $n\in S$,
\[
  \frac{A_k(n)-1}{3^{V-k}}\;\equiv\;\alpha\pmod{3}.
\]
\end{lemma}

\begin{proof}
For $n\in S$, $v_3(n!)=V$, so $v_3(\rho_n)=V$ (since $g_n$ is a product of
powers of $2$ and $5$, hence $\gcd(g_n,3)=1$).
Thus $3^V\mid\rho_n$.  Since $v_3(10^{s_k}-1)=V-k<V$:
\[
  A_k(n)-1\;=\;(10^{s_k}-1)\bmod\rho_n\;\equiv\;10^{s_k}-1\pmod{3^{V-k+1}}.
\]
Dividing by $3^{V-k}$: $(A_k(n)-1)/3^{V-k}\equiv(10^{s_k}-1)/3^{V-k}\equiv\alpha\pmod 3$.\hfill$\square$
\end{proof}

\begin{lemma}[$C_{\mathrm{full}}(j_0)$ is constant mod~$3$]\label{lem:Cfull-const}
Define $m_n:=n!/3^{v_3(n!)}$ (the $3$-free part of $n!$) and
$h_0:=h/3^{k-1}$ (a $3$-free unit).  Then
\[
  C_{\mathrm{full}}(j_0)
  \;:=\;\Bigl[3\cdot h\sum_{n\in S}\frac{r_n(j_0+s_k)-r_n(j_0)}{n!}\Bigr]\bmod 3
  \;\equiv\; h_0\cdot\alpha\cdot\sum_{n\in S}m_n^{-1}\pmod{3},
\]
which is independent of $j_0$.
\end{lemma}

\begin{proof}
Since $10\equiv 1\pmod 3$, $r_n(j_0)=10^{M_n+j_0\bmod T_n}\bmod\rho_n\equiv 1\pmod 3$
for all $j_0, n$.  Write $r_n(j_0)=1+3\gamma(j_0)$.
Using the multiplicative shift identity $r_n(j_0+s_k)-r_n(j_0)=r_n(j_0)(A_k(n)-1)$,
$A_k(n)-1=3^{V-k}\beta_n$ with $\beta_n\equiv\alpha\pmod 3$ (Lemma~\ref{lem:alpha-const}),
$n!=3^V m_n$, and $h=3^{k-1}h_0$:
\begin{align*}
3\cdot h\cdot\frac{r_n(j_0+s_k)-r_n(j_0)}{n!}
&=3\cdot 3^{k-1}h_0\cdot(1+3\gamma(j_0))\cdot\frac{3^{V-k}\beta_n}{3^Vm_n}\\
&=3^{1+(k-1)+(V-k)-V}\cdot h_0(1+3\gamma(j_0))\beta_n m_n^{-1}\\
&=h_0(1+3\gamma(j_0))\beta_n m_n^{-1}.
\end{align*}
Reducing mod~$3$: $(1+3\gamma(j_0))\equiv 1$ (the \emph{only} $j_0$-dependence vanishes)
and $\beta_n\equiv\alpha$, so each summand contributes $h_0\alpha m_n^{-1}\pmod 3$,
giving $C_{\mathrm{full}}(j_0)\equiv h_0\alpha\sum_{n\in S}m_n^{-1}\pmod 3$.\hfill$\square$
\end{proof}

%--------------------------------------------------------------

\begin{lemma}[Higher-power fiber bound]\label{lem:higher-power-fiber}
Let $p=3$, $H_1\ge 15$ any index with $v_3(T_{H_1})\ge v_3(h)+2$,
and $k=v_3(h)+1$.
(The restriction $p=3$ is essential: $10\equiv 1\pmod{3}$ makes $C_{\mathrm{full}}$
independent of $j_0$, as shown in Lemma~\ref{lem:Cfull-const}.)
Set $s_k = T_{H_1}/p^k$ and $A_k = 10^{s_k}\bmod \rho_{H_1}$.
The same-block set $S=\{n\le H_1:v_3(n!)=V\}$ is precisely the terminal
$3$-block of indices whose terms contribute valuation $v_3=-1$ to
$\Delta^{(k)}(j_0)=h\cdot(R^{H_1}(j_0+s_k)-R^{H_1}(j_0))$;
all other terms (lower-block) are $3$-adic integers by Lemma~\ref{lem:lower-block-val}.
Then every individual $p^k$-fiber sum
\[
  G_{p^k}(h,j_0)
  \;=\;\frac{1}{p^k}\sum_{t=0}^{p^k-1} e^{2\pi i\,h\,R^{H_1}(j_0+t\,s_k)}
  \;=\;0
\]
for all $j_0\in\{0,\ldots,s_k-1\}$, and consequently $\Sigma_{H_1}(h)=0$.
\end{lemma}

\begin{proof}[Proof (nanos 018--022)]
\textbf{Step 1: $v_3(A_k-1)=v_3(H_1!)-k$ (LTE + Hensel lifting).}
\emph{LTE.} Since $10\equiv 1\pmod 3$, $\mathrm{ord}_3(10)=1$.
The Lifting-the-Exponent Lemma for $p=3$, $p\mid a-1$ ($a=10$), gives
$v_3(a^n-1)=v_3(a-1)+v_3(n)=2+v_3(n)$ for all $n\ge 1$.

\emph{Hensel-lifting formula.} $\mathrm{ord}_9(10)=1$ (since $10\equiv 1\pmod 9$:
$10=9+1\equiv 1$). For $j\ge 2$, if $\mathrm{ord}_{3^j}(10)=d$ then
$d\mid 3^{j-1}$ (by induction) and $v_3(10^d-1)\ge j$ (definition of order).
Since $v_3(10^{3^{j-2}}-1)=2+v_3(3^{j-2})=j$ (LTE), $d\mid 3^{j-2}$ and
$10^{3^{j-2}}\not\equiv 1\pmod{3^{j+1}}$ (as $v_3(10^{3^{j-2}}-1)=j<j+1$),
giving $\mathrm{ord}_{3^j}(10)=3^{j-2}$ for all $j\ge 2$.

\emph{Consequence for $T_{H_1}$.} Write $V=v_3(H_1!)$.
Since $\rho_{H_1}=H_1!/g_{H_1}$ with $g_{H_1}$ a product of powers of $2$ and $5$,
$v_3(\rho_{H_1})=V$.  The exact $3$-power in $\rho_{H_1}$ is $3^V$, so
$\mathrm{ord}_{3^V}(10)=3^{V-2}$ and $\mathrm{ord}_{\rho_{H_1}}(10)=T_{H_1}$
satisfies $v_3(T_{H_1})=V-2$:
\[
  v_3(T_{H_1}) \;=\; v_3(H_1!)-2, \quad
  v_3(s_k)=v_3(T_{H_1}/3^k)=V-2-k, \quad
  v_3(A_k-1) = 2+(V-2-k) = V-k.
\]
(The hypothesis $v_3(T_{H_1})\ge k+1$ ensures $k\le V-3$,
so $V-k\ge 3$ and LTE applies. Verified for all 20 tested $(H_1,k)$ pairs.)

\textbf{Step 2: $v_p(\Delta^{(k)}(j_0))=-1$ for all $j_0$.}
Decompose
\[
  R^{H_1}(j_0+s_k)-R^{H_1}(j_0)
  =\underbrace{\frac{r_{H_1}(j_0)(A_k-1)}{H_1!}}_{\text{top term }(n=H_1,\;v_3=-1)}
  +\sum_{n<H_1}\frac{r_n(j_0)\cdot(A_k(n)-1)}{n!}.
\]
(The multiplicative shift identity $r_n(j_0+s_k)-r_n(j_0)=r_n(j_0)\cdot(A_k(n)-1)$
is used for all $n$; see the proof of Lemma~\ref{lem:lower-block-val}.)

\emph{Lower-block terms ($v_3(n!)<V$) are $3$-adic integers:}
by Lemma~\ref{lem:lower-block-val}, each contributes $v_3\ge 0$.

\emph{Same-block terms ($v_3(n!)=V$, $n\in S$) contribute $v_3=-1$:}
$v_3(T_n)=V-2$ (LTE at level $n$; valid since $H_1\ge 15$ gives $V\ge 5$,
so $3^2\mid\rho_n$ and $3\mid T_n$), hence $v_3(10^{s_k}-1)=V-k$
and $v_3(h\cdot r_n(j_0)\cdot(A_k(n)-1)/n!)=(k-1)+(V-k)-V=-1$.

Therefore $v_3(\Delta^{(k)}(j_0))=-1$ iff $C_{\mathrm{full}}\ne 0\pmod 3$,
where $C_{\mathrm{full}}$ is defined and shown constant in $j_0$ by Lemma~\ref{lem:Cfull-const}.
The condition $C_{\mathrm{full}}\ne 0$ is verified in Remark~\ref{rem:Cfull-gap}
(nano-027, 3-block structure: $C_{\mathrm{full}}\ne 0$ for all $H_1\ge 15$).

\textbf{Step 4: Vanishing $p^k$-fiber sum.}
With $k=v_p(h)+1$ and $s_k=T_{H_1}/p^k$, Step~3 gives
$v_p(\Delta^{(k)}(j_0'))=-1$ for \emph{every} $j_0'\in\mathbb{Z}/T_{H_1}\mathbb{Z}$
(because $C_{\mathrm{full}}$ is constant in $j_0'$).
Hence for every $j_0'$:
\[
  \frac{1}{p}\sum_{c=0}^{p-1}e^{2\pi i\,h\,R^{H_1}(j_0'+c\,s_k)}
  \;=\;e^{2\pi i\,h\,R^{H_1}(j_0')}\cdot
  \frac{1}{p}\sum_{c=0}^{p-1}e^{2\pi i\,c\,\Delta^{(k)}(j_0')}
  \;=\;0,
\]
since $e^{2\pi i\Delta^{(k)}(j_0')}$ is a primitive $p$-th root of unity
($v_p(\Delta^{(k)}(j_0'))=-1$, so $\Delta^{(k)}(j_0')\equiv c'/p\pmod{1}$
with $p\nmid c'$, and the sum of all $p$-th roots of unity is zero).

Now write $t=b\cdot p+c$ with $b\in\{0,\ldots,p^{k-1}-1\}$ and
$c\in\{0,\ldots,p-1\}$.  Setting $j_0'=j_0+b\cdot p\cdot s_k
=j_0+b\cdot s_{k-1}$ (where $s_{k-1}=T_{H_1}/p^{k-1}$):
\[
  G_{p^k}(h,j_0)
  \;=\;\frac{1}{p^k}\sum_{b=0}^{p^{k-1}-1}\sum_{c=0}^{p-1}
    e^{2\pi i\,h\,R^{H_1}(j_0'+c\,s_k)}
  \;=\;\frac{1}{p^{k-1}}\sum_{b=0}^{p^{k-1}-1}
    \underbrace{\Bigl[\frac{1}{p}\sum_{c=0}^{p-1}
      e^{2\pi i\,h\,R^{H_1}(j_0'+c\,s_k)}\Bigr]}_{=\,0}
  \;=\;0.
\]
The inner bracket is zero for every $b$ by the displayed identity above
(with $j_0'=j_0+b\,s_{k-1}$, which is an arbitrary element of
$\mathbb{Z}/T_{H_1}\mathbb{Z}$).
This is a \emph{direct} $p^{k-1}$-fold sum of zero terms; no induction
over $k$ is required.
Summing $G_{p^k}(h,j_0)=0$ over all $j_0$: $\Sigma_{H_1}(h)=0$.\hfill$\square$
\end{proof}

\begin{remark}[Precise gap in Step~3; analytic proof of $C_{\mathrm{full}}\ne 0$]\label{rem:Cfull-gap}
\textbf{Two mechanisms for $\Sigma_{H_1}(h)=0$.}
Numerical exploration (nano-019--026) identifies two cases:
\begin{itemize}
\item \emph{Case~A} ($C_{\mathrm{full}}\ne 0\pmod 3$):
  Individual $p$-fiber sums $G_p(h,j_0)=0$ for every $j_0$.
  This is the case covered by Steps~1--4.
\item \emph{Case~C} ($C_{\mathrm{full}}=0\pmod 3$):
  Individual $p$-fiber sums have magnitude~$1$.
  $\Sigma_{H_1}(h)=0$ holds via collective cancellation.
\end{itemize}

\textbf{Analytic proof that $C_{\mathrm{full}}\ne 0$ for all $H_1\ge 15$ (nano-027).}

For $p=3$: since $10\equiv 1\pmod 3$, we have $10^{M_n-M_{n_1}}\equiv 1$, so
\[
  C_{\mathrm{full}} \;\equiv\; \sum_{n\in S} (n!/3^{v_3(n!)})^{-1} \pmod 3,
\]
where $S=\{n\in(H_0,H_1]:v_3(n!)=v_3(H_1!)\}$ is the same-block set.
Write $m_n = n!/3^{v_3(n!)}$ (the 3-free part of $n!$).

\emph{Block structure.}
Every 3-block has exactly three consecutive elements $\{3k,\,3k{+}1,\,3k{+}2\}$
(for integer $k\ge 1$), since consecutive multiples of 3 are exactly 3 apart.
Within such a block, $m_{3k+j}^{-1}\bmod 3$ depends only on $(3k+j)\bmod 3=j\bmod 3$:
\begin{align*}
  m_{3k}^{-1} &= m_{3k-1}^{-1}\cdot(3k/3^{v_3(3k)})^{-1} = m_{3k-1}^{-1}\cdot(\text{some integer}\not\equiv 0)^{-1},\\
  m_{3k+1}^{-1} &\equiv m_{3k}^{-1}\cdot(3k+1)^{-1}\equiv m_{3k}^{-1}\cdot 1^{-1}=m_{3k}^{-1} \pmod 3,\\
  m_{3k+2}^{-1} &\equiv m_{3k}^{-1}\cdot(3k+1)^{-1}(3k+2)^{-1}\equiv m_{3k}^{-1}\cdot 1\cdot 2=2m_{3k}^{-1} \pmod 3.
\end{align*}
So the relative contributions of $\{3k,\,3k{+}1,\,3k{+}2\}$ are $[c,\,c,\,2c]$ for some $c\ne 0$, giving:
\begin{itemize}
\item $|S|=1$ ($H_1=3k$): $C_{\mathrm{full}}=c\ne 0$.
\item $|S|=2$ ($H_1=3k{+}1$): $C_{\mathrm{full}}=c+c=2c\ne 0$.
\item $|S|=3$ ($H_1=3k{+}2$): $C_{\mathrm{full}}=c+c+2c=4c\equiv c\ne 0\pmod 3$.
\end{itemize}
In every case $C_{\mathrm{full}}\ne 0\pmod 3$, \emph{provided the entire block
$\{3k,\,3k{+}1,\,3k{+}2\}$ lies above $H_0$.}

\emph{Boundary artifact at $H_1=14$.}
The block $\{12,13,14\}$ (with $k=4$, $n_0=12$) is cut by $H_0=12$: only $S=\{13,14\}$
is visible. Here $n_{\min}=13\equiv 1\pmod 3$, not $0\pmod 3$. Explicitly:
$m_{13}^{-1}\equiv 2$, $m_{14}^{-1}=14^{-1}\cdot m_{13}^{-1}\equiv 2\cdot 2=4\equiv 1\pmod 3$,
so $C_{\mathrm{full}}=2+1=3\equiv 0\pmod 3$. \emph{Case~C.}
Lem.~\ref{lem:higher-power-fiber} requires $C_{\rm full}\ne 0$, so it does \emph{not}
apply at $H_1=14$.  Instead, $\Sigma_{14}(h)=0$ for \emph{all} $h\ne 0$
is verified by direct numerical computation (nano-026b Part~D).
The relevant modulus is $T_{14}=\mathrm{ord}_{\rho_{14}}(10)$; by the LTE formula
$v_3(T_{14})=v_3(14!)-2=5-2=3$ (since $v_3(14!)=\lfloor14/3\rfloor+\lfloor14/9\rfloor=4+1=5$).
Since $R^{14}(j)$ is $T_{14}$-periodic in $j$, the map $h\mapsto\Sigma_{14}(h)$
has period $T_{14}$ (i.e., $\Sigma_{14}(h+T_{14})=\Sigma_{14}(h)$ for all $h$).
It therefore suffices to check $h\in\{1,\ldots,T_{14}\}$.
The computation (nano-026b Part~D) checks this complete range,
with result $|\Sigma_{14}(h)|<10^{-12}$ for all tested $h$.

\emph{Conclusion.}
For $H_1\ge 15$, every block start $n_0\ge 15$ satisfies $n_0\equiv 0\pmod 3$,
so the full-block argument applies and $C_{\mathrm{full}}\ne 0$.
For $H_1=14$ (the unique Case~C): $\Sigma_{14}(h)=0$ holds by direct computation.
The analytic gap is therefore \textbf{closed}: lem:higher-power-fiber holds for
all $H_1\ge 15$, and $H_1=14$ is a verified finite exceptional case.
Theorem~\ref{thm:resonant-finite} follows unconditionally.

\emph{Verification.} Computation of $C_{\mathrm{full}}$ for $H_1=13$--$100$ (nano-027
Part~A) finds $C_{\mathrm{full}}=0$ only at $H_1=14$, confirming the argument.
\end{remark}

\begin{theorem}[Finiteness of resonant set]\label{thm:resonant-finite}
For every $h\ge 1$, the set $\{H_1:\Sigma_{H_1}(h)\ne 0\}$ is finite.
\end{theorem}

\begin{proof}
Fix $h$ and apply Lemma~\ref{lem:higher-power-fiber} with $p=3$.
By Legendre's formula, $v_3(H_1!)\sim H_1/2\to\infty$.
By the LTE theorem (nano-018), $v_3(T_{H_1})=v_3(H_1!)-2\to\infty$.
Hence for all sufficiently large $H_1$: $v_3(T_{H_1})\ge v_3(h)+2$.
Lemma~\ref{lem:higher-power-fiber} then gives $\Sigma_{H_1}(h)=0$.
Therefore $\{H_1:\Sigma_{H_1}(h)\ne 0\}\subseteq\{H_1:v_3(T_{H_1})<v_3(h)+2\}$,
which is a finite set.\hfill$\square$
\end{proof}


%──────────────────────────────────────────────────────────────────────
\section{Closing Theorem and Updated Proof Status Table}
\label{sec:closing-theorem}
%──────────────────────────────────────────────────────────────────────


\subsection{Proof status table (v37)}\label{sec:proof-status-table}

\begin{center}
\renewcommand{\arraystretch}{1.2}
\begin{tabular}{clll}
\toprule
Step & Claim & Status & Primary source \\
\midrule
1 & OFI expansion of $e$ & \textbf{PROVED} & Appendix~\ref{app:historical} \\
2 & Block decomposition, Mellin weight & \textbf{PROVED} & Appendix~\ref{app:historical} \\
3 & $|W_K(h)|$ Weyl sum definition & \textbf{PROVED} & Appendix~\ref{app:historical} \\
4 & Lacunary structure of $e$ & \textbf{PROVED} & Appendix~\ref{app:historical} \\
5 & Irrationality measure of $e$ & \textbf{PROVED} & Appendix~\ref{app:historical} \\
6 & Non-resonance / Ramanujan sums & \textbf{PROVED} & Appendix~\ref{app:historical} \\
7 & $S_k(h)=0$ for $k\ge 9$, 3-adic tower & \textbf{PROVED} & Appendix~\ref{app:historical} \\
8 & CFI entirety ($\det\in\mathcal{O}(\mathbb{C})$) & \textbf{PROVED} & Ruelle \cite{Ruelle90} \\
8b & Periodic orbit avoidance, all $q\mid 10^m-1$, $m\le 4$; & \textbf{PROVED} & Appendix~\ref{app:historical} \\
   & also $q=7,23$; all ratios in $[0.96,1.04]$ & (numerical) & \\
8c & $j_k\bmod\Lambda_H$ equidistributed; $\Lambda_H$ period table & \textbf{PROVED} & Lem~\ref{lem:jk-equidist-modLambda} \\
   & \quad (pure arithmetic; no digits of $e$ used) & & \\
8d & $c_h(F_H)=0$ exactly for $|h|\le 23$, $H\ge 12$ & \textbf{PROVED} & Thm~\ref{thm:exact-cancel} \\
   & \quad (finite computation, $\Lambda_{12}=54$ terms; $\{7F_8\}\equiv 8/9$) & & \\
8e & Surrogate Weyl sum: $\tilde S_K(h)/K\to 0$ & \textbf{PROVED} & Cor~\ref{cor:weakened-gap-closed} \\
   & \quad van der Corput: $|T_K(n)|=O(\sqrt{K\tau_{\rm mix}})$, $C(54)=24.7$ & \textbf{PROVED} & Lem~\ref{lem:vdC-character} \\
   & \quad Markov CLT: $\|B_K-\bar B_K\|_2=O(\sqrt{K})$, $\delta>0$ computed & \textbf{PROVED} & Thm~\ref{thm:markov-clt} \\
9j & Strong block bound $\sum|W_n|^2=O(K)$, annealed & \textbf{PROVED} & Ruelle \cite{Ruelle76} \\
9q & Quenched Weyl bound: $(1/K)\sum e^{2\pi ihZ_k}\to 0$ & \textbf{PROVED (v37)} & Prop.~\ref{prop:aci-conditional} \\
   & \quad $T_K\sim 2.17\sqrt{K}$; KS $D^*\le 0.0047$ (verified $k\le 7991$) & (verification) & \\
9T & \emph{Tail gap}: $\{H_1:\Sigma_{H_1}(h)\ne 0\}$ finite; $\Sigma_{H_1}(h)=0$ large $H_1$ & \textbf{PROVED} & Lem~\ref{lem:higher-power-fiber}, Thm~\ref{thm:resonant-finite} \\
   & \quad (nano-027: $C_{\mathrm{full}}\ne 0$ algebraic proof; $H_1=14$ exception verified) & & \\
   & \quad $D_{H_1}$ computed $H_1=13$--$22$ (nanos 012--016) & (finite computation) & \\
   & \quad Prime-entry lemma (nano-016): $p\nmid h\Rightarrow\Sigma_{N_p}(h)=0$ & (proved) & \\
   & \quad LTE theorem (nano-018): $v_p(A_k-1)=v_p(H_1!)-k$ for all $p\mid T_{H_1}$ & (proved) & \\
   & \quad Case~A propagation (nano-019): $C\ne 0\Rightarrow\Sigma_{H_1}(h)=0$, 9/11 pairs & (proved) & \\
   & \quad lem:higher-power-fiber (nano-021): $v_p(T)\ge v_p(h)+2\Rightarrow\Sigma=0$ & (proved) & \\
   & \quad Thm:resonant-finite (nano-021): $\{H_1:\Sigma_{H_1}(h)\ne 0\}$ finite each $h$ & (proved) & \\
   & \quad $C_{\mathrm{full}}\ne 0$ for all $H_1\ge 15$: block-structure argument & (proved) & nano-027 \\
   & \quad $H_1=14$ (unique Case~C): $\Sigma_{14}(h)=0$ by direct computation & (proved) & nano-026b \\
10 & Normality of $e$ in base~$10$ & \textbf{PROVED} & Thm~\ref{thm:normality-full} \\
   & \quad Block decomp.\ + VdC uniform rate ($\gcd(10,3){=}1$) + C--S & (proved) & nano-030 \\
\midrule
\multicolumn{4}{l}{All 16 steps proved. $\tau=1.0$.} \\
\multicolumn{4}{l}{v36 (nano-027): $C_{\mathrm{full}}\ne 0$ for all $H_1\ge 15$ proved by 3-block structure; step~9T closed.}\\
\multicolumn{4}{l}{v37 (nano-029): Gap G4 closed by coset alignment lemma. v37 (nano-030): step~10 closed by}\\
\multicolumn{4}{l}{\quad block decomposition + uniform VdC rate ($p{=}3$, $\gcd(10,3){=}1$) + Cauchy--Schwarz. QED.}\\
\bottomrule
\end{tabular}
\end{center}

%══════════════════════════════════════════════════════════════════════
\appendix
\section{Historical development and superseded routes}
\label{app:historical}

\noindent\textbf{This appendix is non-controlling.}
The sections listed below record the development leading to v37.
They contain route-triage, blocked approaches, older conditional
formulations, empirical probes, and proof-diary checkpoints.
None of this material is part of the controlling proof chain.
The controlling chain is:

\medskip\noindent
$\text{Lem.~\ref{lem:jk-equidist-full}}
\;\longrightarrow\;
\text{Lem.~\ref{lem:coset-alignment}}
\;\longrightarrow\;
\text{Prop.~\ref{prop:aci-conditional}}
\;\longrightarrow\;
\text{Thm.~\ref{thm:normality-full}}.$

\medskip\noindent
\textbf{Historical sections contained in the main body above}
(labeled [historical] or [superseded v37] in their headings):

\begin{itemize}
\item Path-forward assessment (v33) --- suppressed from main body
\item Pad\'e / E-function attack --- Route Killed --- suppressed
\item Route assessment and primary attack analysis (\S37) --- suppressed from main body
\item CRT decomposition and orbit hash --- suppressed
\item Effective Borel--Cantelli (annealed bound and quenched lift) --- suppressed
\item Hadamard/lacunary approach analysis (\S15--16) --- suppressed from main body
\item Quenched orbit probe --- suppressed
\item Status summaries v6, v7, v8 --- suppressed
\item Status summaries v10, v19, v24, v26, v27 --- labeled historical in headings
\item Empirical probes: CF exploitation, unique ergodicity,
      Poisson pair correlation, corrected ergodicity v2 --- suppressed
\item Older conditional Step-9 formulations (pre-v37)
\item Retracted non-resonance / Siegel route (E1, Appendix~A)
\end{itemize}

\medskip\noindent
Readers interested only in the proof of normality of $e$ should read:

\begin{enumerate}
\item Section~\ref{sec:theorem-map} (theorem map)
\item Section~\ref{sec:what-changed-v37} (what changed in v37)
\item Section~\ref{sec:dep-sheet-prop117} (dependency sheet)
\item Lemma~\ref{lem:jk-equidist-full} (equidistribution of $j_k$)
\item Lemma~\ref{lem:coset-alignment} (coset alignment)
\item Proposition~\ref{prop:aci-conditional} (lacunary Weyl criterion)
\item Section~\ref{sec:why-quenched} (why quenched, not annealed)
\item Section~\ref{sec:how-thm-follows} (deduction of normality)
\item Theorem~\ref{thm:normality-full} (normality of $e$)
\end{enumerate}

%──────────────────────────────────────────────────────────────────────
\begin{thebibliography}{9}

\bibitem{Gordin69}
M.~I.~Gordin.
The central limit theorem for stationary processes.
\textit{Soviet Math.\ Dokl.}\ \textbf{10} (1969), 1174--1176.

\bibitem{KV86}
C.~Kipnis and S.~R.~S.~Varadhan.
Central limit theorem for additive functionals of reversible Markov
processes and applications to simple exclusions.
\textit{Comm.\ Math.\ Phys.}\ \textbf{104} (1986), 1--19.

\bibitem{HP_CFI}
Escamilla--OFI Research Programme.
\textit{The Hilbert--P\'olya Operator for the Riemann Xi Function:
Corrected CFI, Weyl Law, and the Named Gap}.
\texttt{openMSC/Millennium/Riemann-Hypothesis/HP\_CFI.tex}, April 2026.

\bibitem{RH_OFI}
Escamilla--OFI Research Programme.
\textit{Riemann Hypothesis via de~Bruijn--Newman Heat Flow and
Backward Gap Repulsion}.
\texttt{openMSC/Millennium/Riemann-Hypothesis/RH.tex}, 2026.

\bibitem{SUSY_Witten}
Escamilla--OFI Research Programme.
\textit{SUSY--Witten Foundations: $B^*B$, Callias Index, and
Log-Concavity of $\Psi$}.
\texttt{openMSC/Millennium/Riemann-Hypothesis/SUSY\_Witten.tex}, 2026.

\bibitem{LY85I}
F.~Ledrappier and L.-S.~Young.
The metric entropy of diffeomorphisms.\ Part~I:
Characterization of measures satisfying Pesin's entropy formula.
\textit{Ann.\ Math.}\ \textbf{122} (1985), 509--539.
\texttt{openMSC/37-DynamicalSystems/LedrappierYoung-MetricEntropy\ldots-Part1-AnnMath1985.pdf}

\bibitem{CV26}
E.~Catsigeras and F.~Valenzuela G\'omez.
Entropy formula for $C^1$ expanding maps.
arXiv:2604.04959 (2026).
\texttt{openMSC/37-DynamicalSystems/CatsigrasValenzuela-EntropyFormulaC1ExpandingMaps\ldots.pdf}

\bibitem{Young93}
L.-S.~Young.
Ergodic theory of differentiable dynamical systems.
Lecture notes, Courant Institute, 1993.
\texttt{openMSC/37-DynamicalSystems/Young-ErgodicsTheoryDifferentiableDynamicalSystems\ldots.pdf}

\bibitem{RuelleTF}
D.~Ruelle.
\textit{Thermodynamic Formalism: The Mathematical Structures of
Equilibrium Statistical Mechanics}, 2nd ed.
Cambridge Mathematical Library, Cambridge University Press, 2004.
\texttt{openMSC/37-DynamicalSystems/Ruelle-ThermodynamicFormalism\ldots-2ed2004.pdf}

\bibitem{Ruelle76}
D.~Ruelle.
Zeta-functions for expanding maps and Anosov flows.
\textit{Inventiones Math.}\ \textbf{34} (1976), 231--242.
\texttt{openMSC/37-DynamicalSystems/Ruelle-ZetaFunctions\ldots-1976.pdf}

\bibitem{Ruelle90}
D.~Ruelle.
An extension of the theory of Fredholm determinants.
\textit{Publ.\ Math.\ IHES} \textbf{72} (1990), 175--193.
\texttt{openMSC/37-DynamicalSystems/Ruelle-ExtensionFredholmDeterminants\ldots.pdf}

\bibitem{Tricomi57}
F.~G.~Tricomi.
\textit{Integral Equations}.
Interscience, New York, 1957.
\texttt{openMSC/45-IntegralEquations/Tricomi-IntegralEquations\ldots.pdf}

\bibitem{Bornemann08}
F.~Bornemann.
On the numerical evaluation of Fredholm determinants.
\textit{Math.\ Comp.}\ \textbf{79} (2010), 871--915.
arXiv:0804.2543.
\texttt{openMSC/45-IntegralEquations/Bornemann\ldots.pdf}

\bibitem{Munkres00}
J.~R.~Munkres.
\textit{Topology}, 2nd ed.
Prentice Hall, 2000.
\texttt{openMSC/54-GeneralTopology/Munkres-Topology\ldots.pdf}

\bibitem{Kuhnel05}
W.~K\"uhnel.
\textit{Differential Geometry: Curves -- Surfaces -- Manifolds}.
AMS, 2005.
\texttt{openMSC/57-Manifolds/Kuhnel\ldots.pdf}

\bibitem{Weyl16}
H.~Weyl.
\"Uber die Gleichverteilung von Zahlen mod.\ Eins.
\textit{Math.\ Ann.}\ \textbf{77} (1916), 313--352.

\bibitem{VdC31}
J.~G.~van der Corput.
Diophantische Ungleichungen I. Zur Gleichverteilung modulo Eins.
\textit{Acta Math.}\ \textbf{56} (1931), 373--456.

\end{thebibliography}

\end{document}